% ============================================================
% DOCUMENT I: THE ONTOLOGICAL ENGINE
% ============================================================

\documentclass[11pt]{article}

\usepackage{fontspec}
\defaultfontfeatures{Ligatures=TeX,Scale=MatchLowercase}
\setmainfont{Noto Serif}
\setsansfont{Noto Sans}
\setmonofont{DejaVu Sans Mono}
\usepackage{microtype}
\microtypesetup{protrusion=true,expansion=false}
\usepackage{setspace}
\setstretch{1.06}
\setlength{\parindent}{0pt}
\setlength{\parskip}{0.58em}
\setlength{\emergencystretch}{2.5em}
\clubpenalty=10000
\widowpenalty=10000

\usepackage[
  a4paper,
  left=1.00in,
  right=1.00in,
  top=0.88in,
  bottom=0.92in,
  headheight=22pt,
  headsep=0.26in,
  footskip=0.34in
]{geometry}

\usepackage{amsmath,amssymb,amsthm,mathtools}
\usepackage{unicode-math}
\setmathfont{XITS Math}

\usepackage{enumitem}
\usepackage{array,booktabs,longtable,tabularx}
\setlist{
  leftmargin=*,
  itemsep=0.22em,
  topsep=0.38em,
  parsep=0pt,
  partopsep=0pt
}
\setlist[enumerate]{labelsep=0.55em}
\setlist[itemize]{labelsep=0.55em}

\usepackage{xcolor}
\definecolor{seriesInk}{HTML}{1F2933}
\definecolor{seriesAccent}{HTML}{8A3B4A}
\definecolor{seriesSteel}{HTML}{40566F}
\definecolor{seriesMuted}{HTML}{66727E}
\definecolor{seriesPaper}{HTML}{F7F4EE}
\definecolor{seriesRule}{HTML}{C9D0D6}

\usepackage{tcolorbox}
\tcbuselibrary{skins,breakable}
\usepackage{tikz}

\usepackage{titlesec}
\titleformat{\section}
  {\normalfont\Large\bfseries\sffamily\color{seriesInk}}
  {\color{seriesAccent}\thesection}
  {0.72em}
  {}
  [\vspace{3pt}{\color{seriesAccent}\titlerule[0.85pt]}]
\titlespacing*{\section}{0pt}{3.0ex plus .8ex minus .2ex}{1.55ex plus .2ex}

\titleformat{name=\section,numberless}
  {\normalfont\Large\bfseries\sffamily\color{seriesInk}}
  {}
  {0pt}
  {}
  [\vspace{3pt}{\color{seriesAccent}\titlerule[0.85pt]}]
\titlespacing*{name=\section,numberless}{0pt}{3.0ex plus .8ex minus .2ex}{1.55ex plus .2ex}

\titleformat{\subsection}
  {\normalfont\large\bfseries\sffamily\color{seriesSteel}}
  {\color{seriesAccent}\thesubsection}
  {0.68em}
  {}
\titlespacing*{\subsection}{0pt}{2.35ex plus .6ex minus .2ex}{0.82ex plus .1ex}

\titleformat{\subsubsection}
  {\normalfont\normalsize\bfseries\sffamily\color{seriesSteel}}
  {\thesubsubsection}
  {0.62em}
  {}
\titlespacing*{\subsubsection}{0pt}{1.8ex plus .4ex minus .2ex}{0.55ex plus .1ex}

\titleformat{name=\subsection,numberless}
  {\normalfont\large\bfseries\sffamily\color{seriesSteel}}
  {}
  {0pt}
  {}
\titlespacing*{name=\subsection,numberless}{0pt}{2.35ex plus .6ex minus .2ex}{0.82ex plus .1ex}

\renewcommand{\sectionmark}[1]{\markboth{\thesection\enspace #1}{}}
\renewcommand{\subsectionmark}[1]{}

\newtheoremstyle{canonstyle}
  {8pt}{8pt}
  {\normalfont}
  {0pt}
  {\bfseries\sffamily\color{seriesAccent}}
  {.}
  {0.55em}
  {}
\theoremstyle{canonstyle}
\newtheorem{theorem}{Theorem}[section]
\newtheorem{lemma}[theorem]{Lemma}
\newtheorem{proposition}[theorem]{Proposition}
\newtheorem{corollary}[theorem]{Corollary}
\newtheorem{definition}[theorem]{Definition}
\newtheorem{remark}[theorem]{Remark}
\newtheorem{postulate}[theorem]{Physical Postulate}
\newtheorem{principle}[theorem]{Closure Principle}

\providecommand{\Klein}{\mathcal{K}}
\providecommand{\M}{\mathcal{M}}
\providecommand{\R}{\mathbb{R}}
\providecommand{\Z}{\mathbb{Z}}
\providecommand{\Ztwo}{\mathbb{Z}_2}
\providecommand{\Sone}{S^1}
\providecommand{\Sthree}{S^3}
\providecommand{\Diff}{\operatorname{Diff}}
\providecommand{\id}{\mathrm{id}}
\providecommand{\Pin}{\operatorname{Pin}}
\providecommand{\Spin}{\operatorname{Spin}}

\usepackage{fancyhdr}
\usepackage{hyperref}
\usepackage{cleveref}
\usepackage{bookmark}
\usepackage{orcidlink}

\hypersetup{
  colorlinks=true,
  linkcolor=seriesSteel,
  citecolor=seriesAccent,
  urlcolor=seriesAccent,
  filecolor=seriesSteel,
  pdftitle={Shape of Reality: The Necessity of the Zero-Energy Klein Block},
  pdfauthor={Canon},
  pdfsubject={Ontological derivation of the Klein Block topology},
  pdfcreator={LuaLaTeX}
}
\urlstyle{same}

\pdfstringdefDisableCommands{%
  \def\ensuremath#1{#1}%
  \def\mathrm#1{#1}%
  \def\mathbf#1{#1}%
  \def\mathbb#1{#1}%
  \def\mathcal#1{#1}%
  \def\eta{eta}%
  \def\not{not}%
}

\crefname{theorem}{Theorem}{Theorems}
\crefname{lemma}{Lemma}{Lemmas}
\crefname{proposition}{Proposition}{Propositions}
\crefname{corollary}{Corollary}{Corollaries}
\crefname{definition}{Definition}{Definitions}
\crefname{remark}{Remark}{Remarks}
\crefname{postulate}{Physical Postulate}{Physical Postulates}
\crefname{principle}{Closure Principle}{Closure Principles}
\crefname{equation}{Equation}{Equations}
\crefname{section}{Section}{Sections}
\crefname{table}{Table}{Tables}

\pagestyle{fancy}
\fancyhf{}
\fancyhead[L]{\footnotesize\sffamily\color{seriesMuted}\CanonShortTitle}
\fancyhead[R]{\footnotesize\sffamily\color{seriesMuted}\CanonStage}
\fancyfoot[L]{}
\fancyfoot[R]{\footnotesize\sffamily\color{seriesMuted}\thepage}
\renewcommand{\headrulewidth}{0.45pt}
\renewcommand{\footrulewidth}{0pt}
\renewcommand{\headrule}{\hbox to\headwidth{\color{seriesRule}\leaders\hrule height \headrulewidth\hfill}}

\fancypagestyle{plain}{%
  \fancyhf{}
  \fancyfoot[L]{}
  \fancyfoot[R]{\footnotesize\sffamily\color{seriesMuted}\thepage}
  \renewcommand{\headrulewidth}{0pt}
}

\fancypagestyle{titlepage}{%
  \fancyhf{}
  \renewcommand{\headrulewidth}{0pt}
  \renewcommand{\footrulewidth}{0pt}
}

\setcounter{tocdepth}{2}
\usepackage{tocloft}
\setlength{\cftbeforesecskip}{0.42em}
\setlength{\cftbeforesubsecskip}{0.08em}
\renewcommand{\cftsecfont}{\sffamily\bfseries\color{seriesInk}}
\renewcommand{\cftsecpagefont}{\sffamily\bfseries\color{seriesInk}}
\renewcommand{\cftsubsecfont}{\sffamily\color{seriesMuted}}
\renewcommand{\cftsubsecpagefont}{\sffamily\color{seriesMuted}}
\renewcommand{\cftsecleader}{\cftdotfill{\cftdotsep}}
\renewcommand{\cfttoctitlefont}{\Large\bfseries\sffamily\color{seriesInk}}
\setlength{\cftaftertoctitleskip}{1em}

\newcommand{\seriesrule}{%
  \par\vspace{3pt}
  {\color{seriesAccent}\rule{\textwidth}{0.85pt}}
  \par\vspace{2pt}
}
\newcommand{\seriesmark}{%
  \begin{tikzpicture}[baseline=-0.65ex]
    \draw[seriesAccent,line width=0.7pt] (0,0) -- (0.42,0);
    \fill[seriesInk] (0.54,0) circle (1.35pt);
    \draw[seriesAccent,line width=0.7pt] (0.66,0) -- (1.08,0);
  \end{tikzpicture}%
}
\newcommand{\architecturalmark}{%
  \vspace{0.25cm}
  \begin{center}
    {\sffamily\footnotesize\bfseries\color{seriesMuted}
      THREE-PAPER ARCHITECTURE\quad\seriesmark\quad\CanonStage}
  \end{center}
  \vspace{0.15cm}
}

\newtcolorbox{reflect}[1][]{
  enhanced,
  breakable,
  colback=white,
  colbacktitle=seriesPaper,
  colframe=seriesAccent,
  boxrule=0.45pt,
  leftrule=2.4pt,
  arc=1.5pt,
  left=12pt,
  right=12pt,
  top=8pt,
  bottom=8pt,
  before skip=10pt,
  after skip=10pt,
  fonttitle=\bfseries\sffamily\small,
  coltitle=seriesAccent,
  colupper=seriesInk,
  title={#1}
}
\newtcolorbox{keyresult}{
  enhanced,breakable,
  colback=white,
  colframe=seriesInk,
  boxrule=0.75pt,
  arc=1.2pt,
  left=12pt,right=12pt,top=9pt,bottom=9pt,
  before skip=12pt,after skip=12pt,
  fontupper=\bfseries\sffamily\small\color{seriesInk}
}
\newtcolorbox{chainbox}{
  enhanced,breakable,
  colback=seriesInk!4!white,
  colframe=seriesSteel,
  boxrule=0.55pt,
  arc=1.5pt,
  left=14pt,right=14pt,top=10pt,bottom=10pt,
  before skip=12pt,after skip=12pt
}
\newtcolorbox{keyterms}{
  enhanced,breakable,
  colback=seriesPaper,
  colframe=seriesSteel,
  boxrule=0.55pt,
  arc=1.2pt,
  left=13pt,right=13pt,top=9pt,bottom=9pt,
  before skip=12pt,after skip=12pt
}
\newcommand{\ornament}{%
  \par\vspace{5pt}
  \begin{center}\seriesmark\end{center}
  \vspace{1pt}\par
}

\newtcolorbox{plainwords}{
  enhanced,breakable,
  colback=white,
  colframe=seriesMuted,
  boxrule=0.45pt,
  leftrule=2.4pt,
  arc=1.5pt,
  left=12pt, right=12pt, top=8pt, bottom=8pt,
  before skip=10pt, after skip=10pt,
  fonttitle=\bfseries\sffamily\small,
  coltitle=white,
  title={In plain words},
  colupper=seriesInk
}
\newtcolorbox{ifyoudisagree}{
  enhanced,breakable,
  colback=white,
  colframe=seriesMuted,
  boxrule=0.45pt,
  leftrule=2.4pt,
  arc=1.5pt,
  left=12pt, right=12pt, top=8pt, bottom=8pt,
  before skip=10pt, after skip=10pt,
  fonttitle=\bfseries\sffamily\small,
  coltitle=white,
  title={If you disagree},
  colupper=seriesInk
}
\newtcolorbox{forthecurious}{
  enhanced,breakable,
  colback=white,
  colframe=seriesMuted,
  boxrule=0.45pt,
  leftrule=2.4pt,
  arc=1.5pt,
  left=12pt, right=12pt, top=8pt, bottom=8pt,
  before skip=10pt, after skip=10pt,
  fonttitle=\bfseries\sffamily\small,
  coltitle=white,
  title={For the curious},
  colupper=seriesInk
}
\newtcolorbox{howtoread}{
  enhanced,breakable,
  colback=white,
  colframe=seriesMuted,
  boxrule=0.45pt,
  leftrule=2.4pt,
  arc=1.5pt,
  left=12pt, right=12pt, top=8pt, bottom=8pt,
  before skip=10pt, after skip=10pt,
  fonttitle=\bfseries\sffamily\small,
  coltitle=white,
  title={How to read this ledger},
  colupper=seriesInk
}

\newcommand{\CanonTitlePage}{%
\begin{titlepage}
\thispagestyle{titlepage}
\vspace*{0.62cm}
{\sffamily\footnotesize\bfseries\color{seriesMuted}
THREE-PAPER ARCHITECTURE \hfill \CanonStage\par}
\vspace{0.42cm}
{\color{seriesAccent}\rule{\textwidth}{1.05pt}}\par
\vspace{1.65cm}
\ifx\CanonGenre\empty\else
{\sffamily\footnotesize\bfseries\color{seriesMuted}\MakeUppercase{\CanonGenre}}\par
\vspace{0.45cm}\fi
{\sffamily\bfseries\fontsize{28}{32}\selectfont\color{seriesInk}\hyphenpenalty=10000\exhyphenpenalty=10000\CanonTitle\par}
\vspace{0.52cm}
{\sffamily\large\color{seriesAccent}\CanonSubtitle\par}
\vspace{0.78cm}
{\color{seriesAccent}\rule{0.24\textwidth}{0.9pt}}\par
\vspace{0.95cm}
{\normalsize\sffamily Canon\textsuperscript{*}\;\orcidlink{0009-0001-9037-6209}\par}
\vspace{0.25cm}
{\small\sffamily\color{seriesMuted}Independent Researcher\par}
\vspace{0.16cm}
{\small\sffamily\color{seriesMuted}\textsuperscript{*}\href{mailto:canon@necessaryuniverse.com}{canon@necessaryuniverse.com}\par}
\vfill
\ifx\CanonDeck\empty\else
\begin{tcolorbox}[
  enhanced,
  colback=seriesPaper,
  colframe=seriesRule,
  boxrule=0.45pt,
  leftrule=2.4pt,
  arc=1.5pt,
  left=11pt,right=11pt,top=9pt,bottom=9pt,
  width=0.94\textwidth
]
\small\sffamily\color{seriesMuted}
\CanonDeck
\end{tcolorbox}
\vspace{0.65cm}
\fi

\end{titlepage}
}

\newenvironment{CanonAbstract}{%
  \begin{center}
  {\sffamily\footnotesize\bfseries\color{seriesMuted}\MakeUppercase{Abstract}}\\[-0.05em]
  {\color{seriesAccent}\rule{0.16\textwidth}{0.75pt}}
  \end{center}
  \begin{tcolorbox}[
    enhanced,
    breakable,
    colback=white,
    colframe=seriesRule,
    boxrule=0.4pt,
    leftrule=2.2pt,
    arc=1pt,
    left=12pt,right=12pt,top=9pt,bottom=9pt,
    before skip=5pt,after skip=4pt
  ]
  \small
}{%
  \end{tcolorbox}
}

\newcommand{\CanonBibliographySetup}{%
  \small
  \setlength{\parskip}{0.22em}
  \setlength{\itemsep}{0.55em}
}

\newcommand{\CanonStage}{DOCUMENT I \symbol{"00B7} WHY}
\newcommand{\CanonGenre}{Philosophy of Physics}
\newcommand{\CanonTitle}{Shape of Reality}
\newcommand{\CanonSubtitle}{The Necessity of the Zero-Energy Klein Block}
\newcommand{\CanonShortTitle}{Shape of Reality}
\newcommand{\CanonDeck}{A philosophical necessity argument, not a physical experiment or mathematical proof. Every assumption is declared. Document~II supplies the mathematical uniqueness proof; Document~III tests the physical realization. The reader is invited to identify the exact premise or step required for any alternative.}

\begin{document}

\CanonTitlePage

\pagenumbering{roman}
\setcounter{page}{1}

\begin{CanonAbstract}
Physics can describe reality with many mathematical models. This paper asks a more basic philosophical question: \emph{what must be true of reality itself if fundamental reality is to remain logically coherent, non-brute, and physically actualized?} We begin with three declared foundational premises---Identity, the prohibition of Brute Facts, and Physical Actualization---while retaining ordinary logical non-contradiction as the background logical condition of meaningful assertion. The key method is the \textbf{Topological Razor}: if a physical feature is real, it requires a sufficient ground. At every fork, the argument asks: \emph{Can the alternative preserve the declared foundational premises and survive the later admissibility conditions?} If not, it is a rejection of the framework's governing requirements, not a counterexample within it. The five declared bridge conditions are handed to Document~II as explicit hypotheses, and Document~II supplies the mathematical uniqueness proof for the orientation-reversing $S^3$-bundle over $S^1$, called the \textbf{Klein Block}. Zero global energy is then selected separately by the framework's sign-symmetry principle. Consciousness and quantum measurement are treated as interpretive extensions, not load-bearing geometric premises. The result is a conditional necessity argument: within the declared ontology and selection rules, reality cannot be absolute nothingness, and the Klein Block is the uniquely surviving smooth bundle class within the stated category once the declared requirements are retained.
\end{CanonAbstract}

\section*{Architectural Context: The Three-Paper Architecture}
\phantomsection
\addcontentsline{toc}{section}{Architectural Context: The Three-Paper Architecture}

This research program is deliberately divided into three different jobs. Keeping those jobs separate prevents the reader from confusing a philosophical necessity argument with a mathematical uniqueness theorem or a physical realization.

\begin{enumerate}
    \item \textbf{Document~I: The Ontological Engine (This Paper).} \textit{Shape of Reality: The Necessity of the Zero-Energy Klein Block.} This paper establishes the \textbf{WHY}: it states the three foundational premises, develops the elimination rules, and identifies the five geometric conditions that a surviving realization must satisfy.
    \item \textbf{Document~II: The Topological Engine.} \textit{The Klein Block: Topological Uniqueness of the Non-Orientable $S^3$-Bundle over $S^1$}~\cite{CanonWhat}. This paper establishes the \textbf{WHAT}: it takes the five geometric conditions as hypotheses and proves the corresponding smooth bundle classification. Under those hypotheses, the orientation-reversing class is unique and is the Klein Block ($\Klein$). This theorem closes the mathematical escape route; it does not prove that the five hypotheses are physically or philosophically mandatory.
    \item \textbf{Document~III: The Dynamical Engine.} \textit{Twisted Dynamics on the Klein Block: A Conditional Framework for Kinematic Descent, Nodal Defect Structure, and Anomaly Constraints}~\cite{CanonHow}. This paper formulates the \textbf{HOW}: starting from the fixed topology, it develops the conditional dynamical framework and identifies the calculations needed for particle physics, anomalies, baryogenesis, thermodynamics, and other physical tests. It does not complete the Born rule or quantum measurement.
\end{enumerate}

\begin{keyresult}
\textbf{The three-paper logic is therefore:} Document~I eliminates alternatives as admissible physical realizations. Document~II proves the surviving topology is unique within the stated mathematical class. Document~III asks whether that framework-selected structure can realize the physics we observe.
\end{keyresult}

\newpage
\phantomsection
\tableofcontents

\newpage
\pagenumbering{arabic}
\setcounter{page}{1}

\section*{Preface---The Illusion of ``Stuff''}
\phantomsection
\addcontentsline{toc}{section}{Preface---The Illusion of ``Stuff''}

Before we begin the philosophical argument, we need to clear away one deep habit of mind. Nearly everyone carries it: the habit of confusing the user interface with the underlying code.

\begin{reflect}[Question 1---What is real beneath appearances?]
Press your hand against a solid wooden table. It feels hard and unyielding. But what, physically, is actually happening between your hand and the table?
\end{reflect}

At the microscopic scale, your atoms never actually touch the table's atoms. What we call ``touch'' and ``hardness'' is the electromagnetic repulsion between the electron clouds of your hand and the wood. The everyday experience of hardness is not a fundamental property of solid stuff; it is your brain's translation of an invisible force field. 

To understand this, look at your computer screen. You see an icon shaped like a paper folder. You click it, and a window opens. Is there an actual, physical paper folder inside your computer? Of course not. The icon is a \textbf{user interface}. It is a picture the computer draws to help you interact with microscopic electrical states and lines of code. 

Human perception works the exact same way. Your brain is the monitor. Your senses are the interface. 

Consider the world around you:
\begin{itemize}
    \item \textbf{Color:} A photon does not carry ``redness.'' It carries a specific wavelength of electromagnetic energy (say, 650 nanometers). ``Red'' is simply the biological translation your visual cortex renders when that wavelength hits your retina. If there is no eye and no brain to translate it, the wavelength exists, but the color red does not.
    \item \textbf{Temperature:} Molecules possess kinetic energy and momentum. They do not possess ``hotness'' or ``coldness.'' Heat is the biological translation of rapid energy transfer to your nerves.
    \item \textbf{Sound:} Air possesses pressure gradients and mechanical waves. It does not possess ``sound.'' Sound is the auditory translation of those pressure variations hitting an eardrum.
    \item \textbf{Shape:} The visual and tactile experience of a ``cube'' or a ``sphere'' is a biological rendering. The underlying reality is not made of colored, hard shapes; it is made of geometric constraints, topological structures, and field dynamics.
\end{itemize}

Color, temperature, sound, and perceived shape are the universe's desktop icons. They are real, reliable translations, but they are \emph{not} the raw, underlying reality. They are secondary qualities generated by your biological hardware to help a localized subsystem (you) navigate the physical Whole.

\begin{keyresult}
\textbf{The Code vs. The Screen:} The underlying reality---the actual ``blueprint'' of the universe---is not made of sensory impressions. It is described in terms of relationships, fields, geometry, and rules of interaction. 
\end{keyresult}

At the fundamental level, reality is not a warehouse of solid things. It is physically instantiated information: determinate identities, real distinctions, real relations, and real rules. It is not imaginary. It is the self-consistent physical blueprint from which objects, events, and observers become possible.

We must now remove a common misunderstanding about ``geometry.'' In school, geometry is often taught as imaginary math drawn on a chalkboard. But in fundamental physics, geometry is the language we use to describe the actual physical scaffolding of reality. If the true fabric of the universe is expressed through fields and causal connections, then to ask for its ``shape'' is not to reduce it to a picture. It is to ask about its underlying structural code.

This paper asks what shape the underlying code must have if it is to be a complete, self-consistent physical whole. It uses elementary logic to read the blueprint; the advanced mathematics is deliberately handed to Document~II.

The point is not to deny the reality of your experiences. The point is to stop mistaking the desktop icons for the machine itself. We can now state the rules of the argument.

\newpage
\section*{Assumption Ledger}
\phantomsection
\addcontentsline{toc}{section}{Assumption Ledger}

Before we begin the derivation, we declare every load-bearing commitment. The purpose is not to make the conclusion weaker; it is to make the route to the conclusion impossible to hide. A critic should always be able to point to the exact commitment being rejected.

\begin{howtoread}
This table is a map of the entire argument. It lists commitments that will be formally defined and derived in the sections that follow. You do not need to understand them now; they are placed here so you can see exactly what the framework will eventually commit to.
\end{howtoread}

\begin{center}
\small
\begin{tabularx}{\textwidth}{@{}>{\raggedright\arraybackslash}p{4.8cm} >{\raggedright\arraybackslash}p{3.2cm} X@{}}
\toprule
\textbf{Item} & \textbf{Status} & \textbf{Plain meaning} \\
\midrule
$A = A$ & Premise & A determinate thing is itself. \\
Logical Non-Contradiction & Logical background & A proposition and its negation cannot both be true in the same respect. \\
No Brute Facts & Premise & Nothing fundamental is accepted as ``just because.'' \\
Physical Actualization & Premise & A physical truth must be physically instantiated. \\
Compositional Realism & Commitment & The totality of physical existents is one real Whole. \\
Connected Whole / No True Separation & Commitment & Physical reality is one connected totality, not separated by absolute non-being. \\
Topological Razor & Method & No unsupported physical structure is admitted. \\
No-Escape Principle & Principle & A genuine rival must preserve the same foundational requirements. \\
Structural Burden Rule & Method & Positive global structure requires an independent ground; absence of such structure is the default. \\
Read-Head / Global Closure & Postulate & The Whole physically closes upon itself. \\
Finite Actualization & Selection & Fundamental physical content is finite and determinate. \\
Temporal Minimality & Selection & One temporal dimension is admitted. \\
Dimensional Minimality & Selection & Three spatial dimensions are the minimum sufficient choice. \\
No Ungrounded Global Handedness & Selection & No arbitrary permanent global handedness is admitted. \\
Symmetric Selection & Selection & A sign-neutral global quantity is fixed at its symmetry point. \\
Five Closure Postulates & Bridge conditions & The geometric conditions handed to Document~II. \\
Smooth manifold structure & Math idealization & Smooth geometry is the mathematical model used here. \\
Consciousness as interiority & Interpretation & A non-load-bearing interpretation of physical self-reference. \\
\bottomrule
\end{tabularx}
\end{center}

The ledger distinguishes \emph{what we accept} from \emph{what follows once it is accepted}. The three foundational premises are the standards used to judge the alternatives. The named selection principles then apply those standards to particular structural forks.

\newpage
\section{The Three Premises}

We begin with a question that sounds almost too simple: \emph{what must be true before any further question about reality can even make sense?}

\subsection{Premise~1---Identity and Logical Coherence}

\begin{reflect}[Question 2---Can anything exist without being \emph{something}?]
Look at the chair you are sitting on, your own hands, the page or screen in front of you. Are these things \emph{something}, or are they \emph{nothing}?
\end{reflect}

They are something. They are determinate. The chair is the chair; the hand is the hand. We are using the word ``identity'' in this basic logical sense: whatever exists as a determinate thing must be itself.

The next point is just as important. Identity is not the same statement as non-contradiction. $A=A$ is the law of identity. The claim that we cannot have both $A$ and not-$A$ in the same respect is the law of non-contradiction. The two principles are related, but they should not be silently merged.

Why does that matter? Because a proposed alternative must itself be meaningful enough to be proposed. If someone says, ``There is an existent $A$ that is not itself,'' the symbol $A$ must still refer to something determinate in order for the sentence to have content.

\begin{keyresult}
\textbf{Premise~1 (Identity):} Every determinate existent is self-identical: $A=A$.

Ordinary logical non-contradiction is retained as the background logical condition of meaningful assertion: a proposition and its negation cannot both be true in the same respect. This is not an empirical hypothesis.
\end{keyresult}

\begin{ifyoudisagree}
To reject Identity while still making a determinate assertion about what exists is to use the very condition that the assertion denies. The rejection uses what it denies. Unlike the next two premises, Premise~1 is therefore treated as a logical condition of meaningful assertion.
\end{ifyoudisagree}

\subsection{Premise~2---No Brute Facts}

\begin{reflect}[Question 3---Must fundamental features have reasons?]
Ask yourself a sharper question. Suppose we reached the deepest layer of reality. At that final layer, are we entitled to say, ``It is this way, and there is no reason why''?
\end{reflect}

This is the ancient Principle of Sufficient Reason in the form needed here: a demand that fundamental facts not simply stop the search for a reason. Contemporary discussions distinguish such grounding or sufficient-reason claims from ordinary efficient causes; the distinction matters here because the reason for a feature can be structural rather than a prior event~\cite{SEPPSR,SEPGrounding}.

Consider a stone arch. A stone is where it is because of the arrangement and forces of the structure. That is an example of a local fact receiving a reason from a larger structure. Now imagine one stone that has exactly the same surrounding conditions as the others but is simply declared to hover for no reason. That extra fact would be a \textbf{Brute Fact}: a fundamental feature with no sufficient reason.

\begin{keyresult}
\textbf{Premise~2:} Every fundamental feature of reality must have a sufficient reason. Nothing fundamental is accepted as ``just because.'' A Brute Fact is a fundamental feature for which no sufficient ground is available.
\end{keyresult}

Premise~2 is \textbf{independent} of Premise~1. Identity does not imply sufficient reason. We state No Brute Facts as a separate foundational demand. It is strong, controversial, and deliberate. The paper does not pretend otherwise.

\begin{ifyoudisagree}
You are free to reject Premise~2. But then the alternative position has a definite cost: it permits at least one fundamental fact to terminate explanation. That may be a coherent philosophy, but it abandons the no-brute-facts ontology. The disagreement is therefore not hidden. It is located at the foundation.
\end{ifyoudisagree}

\subsection{Premise~3---Physical Actualization}

\begin{reflect}[Question 4---What makes a physical truth physically real?]
A vault is designed to be locked. The blueprint says ``locked.'' The mechanism is capable of locking. But the bolt has never engaged. Is the vault physically locked?
\end{reflect}

No. The blueprint describes a possibility or design. The physical state requires a physical condition that makes the proposition true.

This gives the distinction needed by the present argument. Mathematics can describe a possible structure without asserting that the structure is physically instantiated. A physical truth is different: it is a truth \emph{about the actual physical domain}. The framework therefore requires a physical truthmaker---a physical state, relation, field, geometric condition, or dynamical fact that makes the physical proposition true.

\begin{keyresult}
\textbf{Premise~3:} A truth about physical reality must be physically instantiated. A mathematical description can be true as mathematics without being physically realized; it becomes a physical fact only when the physical domain contains the corresponding truthmaker.
\end{keyresult}

This premise is narrower than a claim about all truths. It does not say that logic or mathematics need physical enforcement. It says only that the adjective ``physical'' carries an actuality requirement.

\begin{ifyoudisagree}
Rejecting Premise~3 allows a theory to call an uninstantiated mathematical possibility a physical fact. That is a coherent change of ontology, but it changes what ``physical reality'' means. The framework does not accept that redefinition.
\end{ifyoudisagree}

These three premises are now in place. The next task is not to guess the universe's shape. It is to ask what alternatives survive all three at once.

\newpage
\section{The Maximal Whole}

We now apply the three premises to the question beneath every cosmology: \emph{could reality fail to exist at all?}

\subsection{The Impossibility of Absolute Nothingness}

\begin{reflect}[Question 5---Could absolute nothingness be real?]
Try to imagine a state of pure, absolute nothingness. No object, no field, no space, no time, no event, no physical law instantiated, no physical relation. Not even an empty container. Could that itself be the actual state of reality?
\end{reflect}

There is a crucial distinction here. We can certainly speak about the \emph{concept} of nothingness. But the question is whether Absolute Nothingness can itself be an \emph{actual physical state}.

Premise~3 says that a physical fact must be physically instantiated. Absolute Nothingness contains no physical instantiation at all. It contains no state that can serve as the physical truthmaker for the claim that a physical reality exists in that condition.

Premise~1 gives the second half of the argument. Whatever is physically real must be determinate enough to be something rather than nothing. Absolute Nothingness is not a hidden kind of physical object. By definition, it contains no physical existent whose identity could constitute an actual physical state.

So the statement ``Absolute Nothingness is the actual physical reality'' does not describe a rival physical world. It removes the physical actuality required for the statement to count as a physical description in the first place.

\begin{keyresult}
\textbf{Within the declared ontology, existence is necessary.} Absolute Nothingness is not an admissible physically actual state within the ontology defined by Identity and Physical Actualization. Therefore some determinate physical actuality must exist.
\end{keyresult}

Within the declared framework, reality cannot be nothing. In the framework's terminology, this is the necessity of existence within the declared ontology. This does not yet prove the specific shape of reality. It proves only that physical actuality cannot be absent.

\begin{ifyoudisagree}
To defend Absolute Nothingness as a genuine rival, one must explain how it can count as an actual physical state while containing no physical actuality. If the proposal instead changes the meaning of ``reality'' so that an entirely empty ontology counts as a reality, that is a different foundational ontology and must be defended on its own terms.
\end{ifyoudisagree}

\subsection{The Existence of the Whole}

\begin{reflect}[Question 6---Is the whole of things a real whole?]
Gather every physical existent---every field, particle, process, region, and physical relation---into one totality. Is that totality itself physically real, or is ``the universe'' only a label invented by the mind?
\end{reflect}

Consider a brick wall. It is made of individual bricks, mortar, and steel. Each part is a real, physical thing. When you combine them, the wall itself becomes a real, physical thing. 

The framework applies this compositional intuition to the physical totality. Every individual particle and field is treated as a real physical existent. By an explicit compositional choice, the sum total of those parts is treated as a real physical existent. We identify this totality as the \textbf{Maximal Whole}, and we denote it by the letter $W$.

\begin{keyresult}
The Maximal Whole ($W$) is a real, physical existent with determinate identity.
\end{keyresult}

The move from ``every part is real'' to ``the sum of the parts is real'' is a mereological principle. Because it is load-bearing, we name it explicitly:

\begin{keyresult}
\textbf{Compositional Realism:} The totality of physical existents is treated as one real, physical Whole. The sum of real parts is itself a real physical existent, not merely an abstract idea in a mind.
\end{keyresult}

Compositional Realism is a load-bearing commitment. It is not a theorem of Identity alone. The reason for stating it is that the argument now concerns \emph{the physical totality itself}. If there is no real Whole, then there is no single physical bearer to which the later requirements of global closure, global topology, and global energy can apply.

\begin{ifyoudisagree}
If you reject Compositional Realism, ``the universe'' becomes a mere mental abstraction rather than a physical reality. Consequently, any truth about the totality lacks a physical truthmaker, violating Premise~3, and the totality itself cannot possess a determinate identity or sufficient reason, violating Premises~1 and~2.
\end{ifyoudisagree}

\subsection{No Outside and No True Separation}

\begin{reflect}[Question 7---Is there an outside or creator?]
Imagine that $W$ is a box. What would the box sit inside? If something physical surrounds it, is that thing outside $W$ or part of $W$? And if a creator is said to exist physically, where does that creator belong?
\end{reflect}

If the thing outside $W$ is physically real, it belongs to $W$ by definition. That is why \textbf{No Outside} is not an extra cosmological guess. It follows from the definition of $W$ as the totality of physical existents.

This gives a precise creator dilemma. A \emph{physical} creator cannot stand outside the Whole: if it physically exists, it is part of the Whole. A \emph{nonphysical} creator is a different metaphysical proposal. It is not an external physical cause inside the ontology studied here; it adds a nonphysical domain and therefore changes the ontology.

What about a gap of Absolute Nothingness between physical parts? When you look at the dark expanse between the stars, it appears to be ``nothing.'' But physics tells us this void is not Absolute Nothingness. It has measurable properties: it bends light, it carries fields, and it permits the propagation of forces. Because it has geometric and physical properties, it possesses a determinate identity (Premise~1) and serves as a physical truthmaker (Premise~3). It is a physical existent, not a gap of non-being.

Empty physical space may be nearly free of ordinary matter, but it is a structural fabric. The framework therefore adds the \textbf{Connected Whole} condition, expressed in the earlier framework language as \textbf{No True Separation}: the physically actual domain is one continuous, connected totality rather than disconnected physical islands separated by regions of absolute non-being.

\begin{keyresult}
\textbf{No Outside:} There is no external physical thing beyond $W$.
\\
\textbf{Connected Whole / No True Separation:} The physical domain is one connected Whole.
\end{keyresult}

Within this ontology the physical Whole is \textbf{uncreated}: no physical creator can stand outside it. Its existence is grounded by the first necessity result, while the detailed form of its internal structure will be fixed by the later closure argument.

\begin{plainwords}
\textbf{Why is there something rather than nothing?} Because Absolute Nothingness does not qualify as an actual physical state under the three foundational premises.

\medskip

\textbf{What is outside the physical universe?} No physical thing. Anything physically real is already part of $W$.

\medskip

\textbf{Who created the physical Whole?} No external \emph{physical} creator can exist. A nonphysical creator is a different ontology and must be defended on that basis.

\medskip

\textbf{Is the Whole just a collection?} Not in this framework. Compositional Realism treats it as one real physical Whole.
\end{plainwords}

\newpage
\section{The Topological Razor}

$W$ exists. $W$ is self-identical. There is no external physical context. Premise~2 now becomes the central question for every remaining structural freedom:

\begin{center}
\textbf{If two physical structures are both mathematically possible, what gives us the right to choose one of them?}
\end{center}

\subsection{The Razor Defined}

\begin{reflect}[Question 8---What rule decides which features are admitted?]
Suppose two possible universes satisfy everything we have said so far. One contains an additional structural feature and the other does not. What gives us the right to declare the extra feature fundamental?
\end{reflect}

The answer is the \textbf{Topological Razor}. It is not a preference for what looks pretty, simple, or familiar. It is Premise~2 applied to physical structure.

In plain language, the Topological Razor says: \textbf{do not admit a feature if doing so requires an unforced choice without a reason.}

If a topological feature is physically real, then it is part of the fundamental physical description. Under No Brute Facts, such a feature cannot simply be present ``for no reason.'' Either another declared principle requires it, an independent physical law grounds it, or the framework has no reason to admit it.

\begin{keyresult}
\textbf{The Topological Razor:} A physical structural feature is admissible only when the framework can identify an independently motivated ground for it---a ground supplied by another declared principle or physical law, not by the mere choice of the branch. If two otherwise admissible choices remain and one introduces an additional ungrounded physical distinction, that branch is rejected by the framework.
\end{keyresult}

The Razor is therefore \emph{motivated by} Premise~2 and is the operational rule by which the premise is applied to topology. We do not claim that Premise~2, considered in isolation, is a theorem about manifolds. The claim is that a no-brute-facts ontology must apply its no-brute-facts requirement to physical geometry as well as to physical matter.

This yields the \textbf{Structural Burden Rule}: positive global structure requires an independent ground; absence of such structure is the default. This asymmetry is not aesthetic; it is logical. In differential geometry, a manifold possesses a generic, unreduced structure group. A ``positive global structure'' (such as global orientability) constitutes a \emph{reduction} of that structure group, adding a global invariant that restricts the physical possibilities. Premise~2 demands a sufficient reason for any such restriction. Conversely, the \emph{absence} of a global invariant means the structure group remains unreduced; it adds no new physical constraint and therefore generates no grounding demand. ``Why is there this extra constraint?'' is a well-formed grounding question; ``Why is there no extra constraint?'' is not.

Before turning the Razor on the universe, it helps to see what it does. Suppose a theory of a coin toss says only that the result must be heads or tails. It then announces ``heads'' without giving any reason why heads rather than tails was selected. Heads is not logically impossible. The problem is that the theory has introduced a fact that its own principles do not select. That is exactly the kind of unexplained addition Premise~2 forbids.

\subsection{The No-Escape Principle}

\begin{reflect}[Question 9---What would actually refute the argument?]
Is it enough to imagine another universe? Or would a genuine refutation have to preserve the same foundational requirements?
\end{reflect}

Here is the rule that governs every fork in the rest of the paper.

\begin{keyresult}
\textbf{No-Escape Principle:} An alternative defeats the present argument only if it preserves the three foundational premises and survives every later admissibility condition while producing a different result. If the alternative succeeds only by rejecting a foundational premise or by adding an ungrounded physical feature, then it is a different ontology, not a counterexample within this one.
\end{keyresult}

This principle does not forbid disagreement. It makes disagreement precise. You may reject the framework. But you may not claim to refute it while secretly rejecting its rules. You may reject No Brute Facts. You may reject Physical Actualization. You may reject global closure. You may defend actual infinity or a fundamental extra dimension. But once you do, the paper will state exactly what has been rejected and what explanatory cost follows.

\subsection{What the Razor Does and Does Not Say}

Let us be precise about the burden of proof.

\textbf{The Razor does NOT say:} ``We cannot think of a reason for feature $X$, therefore $X$ is impossible in every conceivable theory.'' That would be an argument from ignorance.

\textbf{The Razor DOES say:} ``Within this framework, feature $X$ requires a declared ground. If no such ground exists, $X$ is not admitted as a fundamental feature.''

The sections ahead apply this rule repeatedly. At each fork, the question is not ``Can I imagine another option?'' Of course one can. The question is:

\textbf{Can the other option preserve the same premises without introducing a new brute physical fact?}

If yes, the paper has a genuine problem and the alternative must be taken seriously. If no, the alternative has escaped only by changing the starting rules.

The Razor is now active. We can turn it on the physical closure relation of $W$.

\newpage
\section{The Read-Head}

$W$ exists. $W$ is one connected physical Whole. There is no external physical context. We now face the central closure question:

\begin{center}
\textbf{If the Whole has nothing physically outside it, how can its global completeness be physically realized?}
\end{center}

\subsection{The Postulate of Global Self-Closure}

\begin{reflect}[Question 10---What physically enforces global closure?]
A local object can be described by the larger physical environment around it. But $W$ has no larger physical environment. If nothing outside the Whole can supply its final boundary condition, what physically realizes the fact that the Whole is complete?
\end{reflect}

The framework answers with one explicit postulate. We do not hide it inside a later theorem.

\begin{keyresult}
\textbf{The Read-Head Postulate:} The physical Whole possesses a physically realized global closure relation. Its complete physical evolution closes onto itself, so that the physical state associated with the end of the closure is identified with the corresponding state at the beginning. The Read-Head is the name given to this internal closure relation.
\end{keyresult}

The Read-Head is \textbf{postulated}. It is not silently derived from Identity or No Brute Facts. Premise~3 motivates the demand for physical realization; Compositional Realism gives that demand a single physical bearer; and the Read-Head states the specific closure condition that the rest of the argument investigates.

This distinction matters. We are not claiming that Identity alone implies a temporal circle. We are claiming:
\begin{center}
Identity + No Brute Facts + Physical Actualization + global physical self-closure\\ $\Longrightarrow$ a constrained family of possible geometries.
\end{center}

The Read-Head is not an observer outside the universe, not a machine in meta-time, and not a moving spotlight. It is the internal physical closure relation by which the Whole enforces its own self-identity.

\subsection{Consciousness and Measurement}

\begin{reflect}[Question 11---What does self-reference have to do with consciousness?]
If the Whole is physically self-closing, is consciousness another substance standing outside it, or could conscious self-reference be an organized physical process inside the same Whole?
\end{reflect}

On this framework's interpretation, consciousness is the \textbf{interiority of physical self-reference}. A human mind is one of the localized, highly organized self-referential subsystems of the same physical Whole whose global closure is described by the Read-Head.

This is an interpretation, not a premise needed to obtain the topology.

\begin{keyresult}
\textbf{The Consciousness Interpretation:} Consciousness and quantum measurement are interpretive extensions of the closure ontology. They are not load-bearing for the five geometric postulates. Rejecting this interpretation leaves the topological argument unchanged.
\end{keyresult}

\subsection{The Read-Head's Requirements}

The Read-Head is a physical closure relation. For it to do the work assigned to it, the framework requires three features:
\begin{itemize}
    \item \textbf{Sequence:} there must be an ordered relation between physical states. This is the role assigned to time.
    \item \textbf{Extension:} physical states must have an arena in which distinct regions and interacting processes can exist. This is the role assigned to space.
    \item \textbf{Completion:} the closure relation must be complete rather than permanently unfinished.
\end{itemize}

\subsection{The Five Closure-Admissible Postulates}

The preceding requirements generate five specific bridge conditions that the framework insists a surviving geometric realization must satisfy.

\begin{enumerate}
    \item \textbf{Global Temporal Fibration.} The Whole admits a smooth fibration (a stacking of space-slices along time) over a connected, compact one-dimensional temporal base without boundary (no edge), with connected spacelike three-dimensional fibers. Compactness is interpreted as completed finite physical extent, not finite mathematical cardinality.
    \item \textbf{Spatial Compactness and Boundarylessness.} Each spatial fiber is compact and without boundary.
    \item \textbf{Spatial Simple Connectivity.} Each spatial fiber is simply connected (no fundamental holes).
    \item \textbf{Time-Orientability.} A consistent local choice of future direction exists throughout spacetime. (Document~II later proves this is mathematically automatic once a Lorentzian metric with spacelike fibers is established, but it is retained here as an explicit physical architectural requirement.)
    \item \textbf{Nontrivial Vertical Orientation Monodromy.} After one circuit of the temporal base (what happens to the spatial slice after one full trip around the time loop), the orientation of the spatial fiber is reversed (no permanent global handedness).
\end{enumerate}

These five conditions are the bridge from philosophy to mathematics. Document~I argues why the framework requires them. Document~II takes them as exact hypotheses and proves what mathematical structure remains.

\begin{plainwords}
The Read-Head is not a hidden observer and not a machine outside reality. It is the name for the internal physical closure relation required by the framework. From here on, every choice must answer the same question: \textbf{what gives this feature its reason?}
\end{plainwords}

\newpage
\section{Finiteness}

We now turn the Topological Razor on the size of the Whole. There are two ideas that must not be confused:
\begin{center}
\textbf{Mathematical infinity} \qquad\text{vs.}\qquad \textbf{Actual physical infinity}.
\end{center}

\subsection{The Read-Head Must Complete}

\begin{reflect}[Question 12---Can a physical closure be permanently unfinished?]
Suppose the Whole has a closure relation whose completion is always one step away. Has the Whole actually closed, or has the closure merely been postponed forever?
\end{reflect}

The Read-Head is postulated as a completed physical closure relation. An endless process that never reaches its stipulated closure does not realize that closure. 

An unfinished closure is not the same thing as an actually infinite object. A physical Whole might be infinite and still be described by a rule. Therefore the argument must address actual physical infinity directly.

\subsection{The Case Against Actual Physical Infinity}

\textbf{Abstract mathematical infinity} is fully legitimate mathematics. Nothing in this paper says otherwise.

\textbf{Actual physical infinity} is different. It says that an infinite physical collection or extent is itself a completed physical fact.

\subsubsection*{The Great Hall Illustration}

\begin{reflect}[Question 13---Does infinity behave like a finite physical amount?]
Imagine a hall with infinitely many occupied rooms. Remove every guest in the odd-numbered rooms. Have you removed an infinite number of guests? Yes. Does the remaining collection still have the same cardinality as before? In the mathematics of countable infinity, yes.
\end{reflect}

This is the familiar Hilbert's Hotel lesson~\cite{SEPInfinity}. It does \emph{not} prove that actual physical infinity is impossible. It shows something narrower and important: mathematical cardinality does not behave like an ordinary finite measure of physical amount. Hilbert's Hotel is not the proof. It illustrates why physical determinacy requires Finite Actualization.

\begin{keyresult}
\textbf{Finite Actualization:} A physically completed Whole must have finite, determinate physical content. Mathematical infinity remains legitimate as mathematics. Actual physical infinity is excluded by this framework because a completed physical Whole is required to have a determinate finite total content.
\end{keyresult}

\begin{ifyoudisagree}
You may defend an actually infinite physical Whole. But then you must reject Finite Actualization or explain why a completed physical totality may possess actual infinity while still satisfying the framework's requirement of finite, determinate physical content.

\medskip

The cost is serious. If you accept actual physical infinity, you cannot also demand that the Whole be physically completed. Total questions---how much physical content, how many physical parts, what final closure completes the Whole---no longer have determinate finite answers unless you supply a new grounding principle.
\end{ifyoudisagree}

\subsection{The Continuum and Physical Content}

Finite Actualization concerns \emph{physical content}, not the number of mathematical coordinate points in a continuum model. A smooth manifold may contain infinitely many points and still represent a finite physical volume. The framework is therefore not confusing mathematical cardinality with physical quantity.

\subsection{Finite but Boundaryless}

\begin{reflect}[Question 14---Must every finite thing have an edge?]
A finite sheet can end at an edge. Must a finite universe therefore end at a wall?
\end{reflect}

No. Finiteness and boundarylessness are different properties. A boundary is additional structure. If the universe ends at a particular physical boundary, Premise~2 asks for the reason that the boundary is there and not somewhere else. The framework has supplied no independent physical ground for such a wall.

The familiar surface of a balloon gives the simple picture. Its total area is finite, but an ant moving along the surface need never encounter an edge. The important distinction is mathematical: \emph{compact} means topologically complete; the framework interprets it as finite physical extent, not finite mathematical cardinality. It does not mean ``has a wall.''

\begin{plainwords}
Mathematical abstract infinity concept is not the problem. The question is whether an \emph{actual physical infinity} can be a completed fundamental fact. This framework says no. And a finite Whole does not need a wall: compactness and boundarylessness can coexist.
\end{plainwords}

\newpage
\section{Time Is a Circle}

We have now fixed the temporal problem at the level of structure: the Whole must have a completed closure, and the Read-Head supplies an ordered sequence within that closure.

\begin{reflect}[Question 15---What shape can a completed sequence have?]
If time is the dimension in which one physical state is ordered relative to another, what happens when the complete physical history itself has no beginning and no end?
\end{reflect}

\subsection{The Dimension of Sequence}

The framework assigns the role of \textbf{time} to the dimension that orders physical sequence. 

\subsection{One Dimension Only}

The Read-Head needs one ordered sequence: one state follows another along a temporal direction. A second time dimension is not required for that function. This gives \textbf{Temporal Minimality}: when one temporal dimension performs the required ordering function, no additional temporal dimension is admitted without an independent reason.

\begin{ifyoudisagree}
Defend a second temporal dimension if you wish. But then identify the physical function it supplies that the first does not. Merely showing that two-time mathematics can be written does not supply the missing physical ground.
\end{ifyoudisagree}

\subsection{No Endpoints and the Only Shape}

Could the temporal sequence have a first moment or a last moment? Within this closure ontology, either endpoint would be a distinguished temporal boundary. Premise~2 therefore asks: \emph{what physically grounds that exact boundary?} No independent ground is supplied. The endpoint branch is excluded.

We are left with a temporal dimension that is one-dimensional, finite, connected, compact, and boundaryless. The standard classification gives exactly one possibility:

\begin{keyresult}
The shape of time is $\Sone$. It is a compact one-dimensional temporal base with no first or last point.
\end{keyresult}

Think of a clock face. It has no privileged first mark built into the circle itself. 

\subsection{Closed Time Is a Real Cost and a Real Test}

A closed temporal base has a serious consequence. In a Lorentzian realization where the temporal base direction is timelike, the closed temporal direction produces closed timelike curves. This is a genuine causal feature of the resulting spacetime and has a substantial literature, including the chronology-protection debate~\cite{HawkingChronology}. The framework therefore does not hide this consequence. It accepts the topological closure and requires the eventual dynamics to show that the resulting causal structure can be physically viable.

\newpage
\begin{forthecurious}
Document~II constructs a time-orientable Lorentzian realization and proves (via the Lefschetz fixed-point theorem) that \emph{every} smooth future-directed timelike vector field transverse to the spatial fibers possesses a closed orbit. Consequently, no global real-valued time function can exist. The topology carries a real, mathematically proven causal price, and its physical viability belongs to the dynamical tests of Document~III.
\end{forthecurious}

\subsection{The Arrow of Time and Entropy}

Time-orientability is a separate requirement. The framework retains \textbf{Time-Orientability}: a continuous local choice of future direction exists throughout the spacetime.

\begin{reflect}[Question 16---If time is a circle, how can entropy keep increasing?]
Suppose an entropy-like quantity strictly increases as one moves toward the future. After one complete circuit of the temporal circle, where has the increase gone?
\end{reflect}

It cannot have increased strictly and returned to its original value. Therefore a single-valued entropy scalar cannot be globally monotone on the entire temporal circle. This is a dynamical obligation. A day can warm and cool, but if time is a loop the temperature curve must eventually reconnect. Entropy cannot simply rise forever around the loop. The Big Bang and Heat Death are treated as coordinates on the loop, not absolute endpoints.

\begin{forthecurious}
Document~II proves a strict topological no-go theorem (Theorem 4.9): a nonconstant, single-valued scalar entropy potential factoring through the temporal base cannot be globally monotone around the circle. Document~III strengthens this (Proposition 3.7) by showing that any globally defined, quotient-compatible entropy current with pointwise nonnegative divergence must have vanishing integrated entropy production. The physical resolution of this Tolman entropy problem (Open Problem 22 in Document~III) requires a non-circular mechanism, such as a thermodynamic algebra transmutation at a topologically forced nodal crossing.
\end{forthecurious}

\begin{plainwords}
Time is a circle because the framework demands a completed, one-dimensional, boundaryless temporal closure. The local future direction can remain consistent. The price is real: global causal and thermodynamic behavior is constrained by no-go results in Documents II and III, but its physical completion remains an explicit dynamical test.
\end{plainwords}

\newpage
\section{Space Is a 3-Sphere}

Time supplies ordered sequence. The Read-Head also requires \textbf{extension}: physical processes need a spatial arena in which distinct regions and interactions can exist.

\begin{reflect}[Question 17---How many spatial dimensions are actually required?]
Could reality have one, two, three, or ten spatial dimensions? What does the Read-Head actually require before extra dimensions become additional ungrounded structure?
\end{reflect}

\subsection{Space Is Physical}

The framework treats spacetime geometry as part of physical ontology. The space between physical systems is part of the geometric structure in which those systems are related, not absolute nothingness.

\subsection{Why Exactly Three Dimensions?}

Imagine building a complex, working machine with interacting parts on a flat, two-dimensional surface. You must lay down physical pathways to connect the parts. In a simple two-dimensional layout, pathways that must cross cannot pass without meeting, interfering, or requiring some additional structure. A third spatial direction gives the pathways room to pass without meeting.

This is a picture to help the mind. It is not the proof.

The deeper point is this: under the framework's interpretation of extension, the spatial arena must be able to host general finite networks of physical relations, modeled as non-intersecting channels, without forced intersections. Two dimensions are not sufficient for arbitrary finite networks of that kind. Three dimensions are sufficient.

This yields \textbf{Dimensional Minimality}: once three dimensions perform the required extension function, additional spatial dimensions require an independent physical ground.

\begin{keyresult}
\textbf{Three spatial dimensions:} Under the framework's stated interpretation of extension, three is the minimum sufficient spatial dimensionality. This is a framework-level minimality result under the stated extension requirement, not a claim that higher-dimensional mathematics is inconsistent.
\end{keyresult}

\subsubsection*{The Cost of More Than Three}

If three dimensions already perform the required extension function, any additional spatial dimension is extra physical structure. Under Premise~2, the number of dimensions is itself a physical fact and cannot be accepted as ``just because.''

Once Dimensional Minimality is abandoned, the framework has no internal reason to stop at four, ten, ten billion, or any larger finite number. The dimensionality becomes an ungrounded brute number, and every higher finite value becomes an equally ungrounded candidate unless a new selector is supplied. An actually infinite number of spatial dimensions is not a solution, because Finite Actualization excludes actual physical infinity.

This is not an infinite regress of dimensions causing dimensions. It is an unbounded brute-choice regress: without a grounding principle, the number of spatial dimensions has no stopping rule.

There is also an empirical cost. The observed macroscopic physical arena appears three-dimensional. If additional fundamental spatial dimensions are admitted, the theory must explain why they are not observed. If they are hidden, compactified, or dynamically suppressed, that hiding mechanism is itself a new physical fact requiring grounding. If they are not hidden, the proposal conflicts with the empirical signal of three-dimensional space.

Therefore, accepting more than three dimensions is not forbidden because higher-dimensional mathematics is impossible. It is inadmissible within this framework unless the extra dimensions are independently grounded, their exact number is selected, and their apparent absence from macroscopic observation is explained.

\begin{forthecurious}
A standard graph-theoretic fact makes the minimality point precise: the complete graph on five points cannot be drawn in a plane without crossings, but it can be embedded in three-dimensional space without forced crossings. More generally, finite networks modeled as non-intersecting channels can be embedded in three dimensions without forced crossings.
\end{forthecurious}

\begin{ifyoudisagree}
You may challenge this step in three ways. First, reject the requirement that general finite physical networks, modeled as non-intersecting channels, must be embeddable without forced intersections. Second, show that the Read-Head's physical function can be realized in fewer than three dimensions under a different interpretation of extension. Third, accept more than three dimensions, but then supply: (i) an independent physical ground for the extra dimensions; (ii) a selector for their exact number; and (iii) an account of why the observed macroscopic arena appears three-dimensional. Without these, the extra dimensions are ungrounded structure.
\end{ifyoudisagree}

\subsection{Finite, Edgeless, and Hole-Free}

The spatial fibers inherit the closure requirements. First, they are selected to be compact and without boundary. A hard spatial wall would require a ground.

Second, what about holes? Imagine a torus rather than a sphere. A loop winding around the torus can fail to shrink to a point. That nontrivial loop structure is an additional topological feature. Extra loop structure would be additional topology without independent ground. The framework accordingly imposes \textbf{Spatial Simple Connectivity}: every closed spatial loop can be continuously shrunk to a point.

\subsection{The Only Possible Shape}

Now ask the mathematical question: What connected, compact, boundaryless, simply connected 3-manifold is left?

\begin{keyresult}
Space is $\Sthree$.
\end{keyresult}

The 3-sphere is compact and boundaryless: finite spatial extent without a wall. In the standard round geometry, straight paths can return to their starting point, but topology alone does not guarantee closed geodesics for every possible metric.

\begin{forthecurious}
The result is the Poincar\'e theorem: every closed simply connected 3-manifold is homeomorphic to $\Sthree$, with the corresponding smooth classification in dimension three supplied by Moise's work~\cite{Perelman,Moise}. Metric behavior belongs to the dynamical realization in Document~III.
\end{forthecurious}

\begin{plainwords}
Space is a finite, boundaryless, simply connected three-dimensional whole. Once those conditions are imposed, the mathematical classification leaves $S^3$ as the unique spatial topology.
\end{plainwords}

\newpage
\section{The Twist}

Space is $\Sthree$. Time is $\Sone$. The last geometric question is decisive:

\begin{center}
\textbf{How must the spatial slices be identified when the temporal circle closes?}
\end{center}

\begin{reflect}[Question 18---How are time and space glued together?]
Take a 3-sphere of space and move once around the temporal circle. When you return to the starting time, must the spatial orientation return unchanged, or can the identification reverse it?
\end{reflect}

\subsection{Why the Dimensions Add}

Sequence requires ordered stages. Each stage must be an extended physical arena. Therefore the Whole is modeled as a stack of spatial slices ordered along the temporal base.

Once this fibration (stacking of space-slices along time) is accepted, the dimension count is immediate:

\begin{center}
Dimension of Whole = dimension of fiber + dimension of base = 3 + 1 = 4.
\end{center}

The four-dimensional character is therefore a mathematical consequence of the chosen fiber and base dimensions together with the fibration structure. The existence of the global fibration itself is a bridge condition of the framework.

\subsection{Stacking the Slices}

Picture a loaf of bread. Each slice is a moment of time (a 3-sphere of space). The whole loaf is the universe. Now, because time is a circle, you must bend the loaf around and glue the first slice to the last. How exactly do you attach them? This is a picture to help the mind. It is not the proof.

Document~II proves that, under the stated smooth bundle hypotheses, there are exactly two smooth $S^3$-bundle classes over $S^1$~\cite{CanonWhat}:
\begin{enumerate}
    \item \textbf{Orientation-preserving gluing:} the spatial orientation returns to itself after one circuit.
    \item \textbf{Orientation-reversing gluing:} the spatial orientation returns reversed after one circuit.
\end{enumerate}

This is a complete mathematical fork within the specified category. These exhaust the smooth bundle classes in the stated category.

\begin{forthecurious}
The mathematical proof that there are exactly two classes relies on Hatcher's proof of the Smale Conjecture~\cite{Hatcher}, which establishes that the diffeomorphism group of the 3-sphere deformation-retracts onto the orthogonal group $O(4)$. This restricts the mapping class group to $\mathbb{Z}_2$, leaving only the orientation-preserving and orientation-reversing components.
\end{forthecurious}

\subsection{The Straight Glue and the Global Handedness Test}

The orientation-preserving branch carries a globally preserved fiber-orientation structure. The framework asks whether that structure is independently grounded.

Does the rest of the framework require that structure? No. The Read-Head requires a temporal ordering. It does not require a globally preserved spatial orientation. The distinction is important: we are not saying that mathematics becomes inconsistent if the bundle is orientation-preserving. We are saying that the orientation-preserving class carries a global structural feature---the preservation of fiber orientation around the loop---that is not grounded by any earlier requirement.

The asymmetry used at this fork must be stated precisely in topological terms. The orientation-preserving branch is \emph{globally orientable}: it admits a continuous, globally consistent choice of spatial orientation around the temporal loop. The orientation-reversing branch is \emph{non-orientable}: it lacks any such globally consistent orientation. The framework therefore applies the \textbf{Structural Burden Rule}: the existence of a global orientation is a positive structural feature that requires an independent ground; the absence of a global orientation is not an additional structure of the same kind. This rule is a declared methodological commitment of the framework, not a theorem of the mathematics.

This is where \textbf{No Ungrounded Global Handedness} enters. It is the Topological Razor applied to this exact two-way mathematical fork:

\begin{keyresult}
\textbf{No Ungrounded Global Handedness:} When the mathematical classification leaves an orientation-preserving and an orientation-reversing class, and no independent principle requires a globally preserved fiber orientation, the orientation-preserving class is rejected as an ungrounded global distinction.
\end{keyresult}

The logic is not ``twist looks simpler.'' It is: No Brute Facts $\Longrightarrow$ no fundamental ungrounded global structural distinction. The straight branch retains a global fiber-orientation structure that the present framework does not independently ground; the twisted branch does not admit such a globally preserved orientation.

\begin{ifyoudisagree}
There are only two ways to preserve the straight branch within the present framework. First, provide an independent physical ground for global preservation of spatial orientation. Second, reject No Ungrounded Global Handedness. Simply saying ``the straight bundle is mathematically possible'' is not enough: Document~II already establishes that. The philosophical question is whether its extra global orientation structure is grounded.
\end{ifyoudisagree}

\subsection{The Twist Is the Surviving Branch}

Take a leather belt. Buckle it normally, and the smooth side is always out. Now unbuckle it, give it a half-twist, and buckle it again. As you trace your finger along it, the smooth side becomes the rough side. The twist prevents a permanent global distinction from being carried unchanged around the loop. This is an analogy, not a literal physical prediction.

The physical geometry is more precise than the analogy. The orientation-reversing monodromy means that spatial orientation cannot be globally chosen so as to return unchanged after one complete circuit of time.

\begin{keyresult}
\textbf{Klein Block:} Within the declared premises, closure conditions, and No-Escape Principle, the orientation-preserving class is excluded unless independently grounded. The unique surviving smooth bundle class is the orientation-reversing $S^3$-bundle over $S^1$, called the \textbf{Klein Block} ($\Klein$)~\cite{CanonWhat}.
\end{keyresult}

This is the point at which the philosophy and mathematics meet most cleanly. The philosophy identifies the unacceptable extra distinction; the mathematical classification guarantees that, once that branch is removed, only one smooth bundle class remains.

\begin{forthecurious}
We call it the \textbf{Klein Block} by analogy with the classical Klein bottle. The historical name ``Klein'' belongs to Felix Klein's work associated with the Klein bottle; the present four-dimensional object is a higher-dimensional analogue of the same orientation-reversing bundle idea~\cite{Klein}.
\end{forthecurious}

\subsection{The Local Arrow of Time Remains}

The twist acts on the spatial fibers. It does not reverse the local future direction of the temporal base.

The framework therefore keeps \textbf{Time-Orientability}: locally, one timelike direction is continuously selected as future. Global closure of time and global reversal of spatial fiber orientation are different statements.

\begin{forthecurious}
Document~I does not by itself prove that the seam is CPT. Any CPT interpretation belongs to the dynamical theory of Document~III. In the wider closure ontology, the seam may be treated as a CPT-type closure only if dynamically established.
\end{forthecurious}

\subsection{The Seam and Open Dynamical Tests}

The circle and the twist do two different jobs. The circle supplies global temporal closure. The twist removes the need for a globally preserved spatial orientation.

Local left and right remain meaningful. What disappears is the possibility of carrying one globally fixed spatial orientation consistently through the entire temporal loop.

The twist precludes a globally preserved spatial orientation from being carried unchanged around the temporal loop. It provides a structural condition relevant to chirality, but it does not by itself derive weak-interaction parity violation or baryogenesis~\cite{LeeYang,Wu,CanonHow}.

As an analogy: if one could follow the full circuit of the loop, the seam identification relates the state to its orientation-reversed counterpart. This is a picture to help the mind, not a literal prediction that an astronaut physically flips.

\begin{forthecurious}
Mathematically, the Klein Block can be written explicitly as a \textit{mapping torus}:
\[
\Klein = \frac{\Sthree \times [0,L]}{(x,0)\sim(r(x),L)},
\]
where $\Sthree$ is the spatial 3-sphere, $[0,L]$ represents one circuit of the temporal circle, and $r$ is the orientation-reversing reflection. This is the exact mathematical object classified by Document~II~\cite{CanonWhat}. 

\medskip

A crucial philosophical caveat: the interval $[0,L]$ should not be mistaken for an ``external time'' in which the universe is manufactured. It is merely a fundamental domain for the circle base. The endpoint identification \emph{is} the closure relation.

\medskip

What does this twist do to the physics? Document~III explores whether this orientation-reversing geometry forces certain physical fields (specifically, CP-odd condensates) to drop to zero at specific topological defect lines (an $S^2 \times S^1$ locus). If it does, this topological ``kink'' could potentially seed the parity violation and matter-antimatter asymmetry (baryogenesis) we observe in the real universe. However, proving that this mechanism actually works remains one of the explicit open dynamical tests of the architecture~\cite{CanonHow}.
\end{forthecurious}

\begin{plainwords}
The mathematics gives two possible ways to glue the spatial $S^3$ around the temporal $S^1$. The framework rejects the straight branch because it carries an ungrounded global orientation distinction. The twisted branch is therefore the unique surviving smooth bundle class: the Klein Block.
\end{plainwords}

\newpage
\section{Zero Energy}

The topology is now fixed within the framework. A separate question remains: what global energy value can the completed Whole possess without introducing another ungrounded choice?

\begin{reflect}[Question 19---What is the total energy of the Whole?]
Suppose the Whole has one global energy value. Why is that value positive, negative, or zero? And what would it mean for one of those values to be selected without a reason?
\end{reflect}

\subsection{The Problem With an Arbitrary Global Number}

The framework does not say that local physical energies vanish. Stars can shine, particles can move, fields can store energy, and gravitational systems can have nonzero local or subsystem energies.

The question is narrower: suppose there is a meaningful global energy datum assigned to the Whole. Call it $E_{\mathrm{global}}$. Premise~2 then asks a simple question: why that value rather than another? The answer cannot simply be ``because that is the value.'' That would make the global value a Brute Fact.

\subsection{Symmetric Selection}

Consider the sign-reversal symmetry: $E_{\mathrm{global}} \mapsto -E_{\mathrm{global}}$.

Every nonzero value selects a sign and a magnitude. Zero selects neither.

Under \textbf{Symmetric Selection}, the framework requires a globally sign-neutral value when no independent principle grounds a preferred sign. A value that is changed by the symmetry cannot be selected without an additional choice. The only fixed point satisfies $E_{\mathrm{global}} = -E_{\mathrm{global}}$, and therefore $E_{\mathrm{global}} = 0$.

This is not a numerical trick. It is a symmetry argument. Once the sign-reversal symmetry is admitted as intrinsic to the global energy selection problem, zero is the unique value that requires neither a sign choice nor a compensating magnitude.

\begin{keyresult}
\textbf{Zero-Energy Closure Selection:} Under the framework's Symmetric Selection principle, the global energy value is fixed at the unique fixed point of $E \mapsto -E$:
\[
E_{\mathrm{global}} = 0.
\]
A nonzero value remains available only by rejecting the symmetry-selection requirement or by supplying an independent physical ground for the sign and magnitude.
\end{keyresult}

Zero total energy is not one of the five topological postulates passed to Document~II. It is a separate closure selection.

\begin{ifyoudisagree}
The real point of disagreement is now visible. You may argue that the sign-reversal symmetry is not an intrinsic symmetry of the global energy concept. That is a direct challenge to Symmetric Selection. What you may not do is accept the symmetry, choose a nonzero value anyway, and then call that choice non-arbitrary. Once the symmetry is part of the framework, zero is its only fixed point.
\end{ifyoudisagree}

\subsection{Zero Energy and the Gravitational Constraint}

The preceding paragraph is the framework's philosophical selection. It should not be confused with a statement that General Relativity possesses an ordinary ADM energy for the compact universe and that the measurement has been found to be zero.

For compact, boundaryless spatial slices, the standard asymptotic ADM energy at spatial infinity is not the relevant global quantity. In the canonical formulation of General Relativity, the Hamiltonian is built from the gravitational and matter constraints, together with boundary terms when appropriate. On a compact boundaryless slice, there is no spatial-infinity boundary term; on the constraint surface the canonical Hamiltonian vanishes weakly~\cite{Wald,ADM}. Document~II states this carefully on the infinite cyclic cover~\cite{CanonWhat}.

The logical relationship is therefore: the framework's global zero selection and $H_{\mathrm{canonical}} \approx 0$ are distinct statements that are mutually compatible in the intended realization. The second does not, by itself, prove the first. The first is the philosophical closure selection; the second is a dynamical consistency result.

\begin{forthecurious}
Document~II deliberately calls the Hamiltonian statement conditional. It is a statement about the constrained canonical Hamiltonian on the lifted compact slices, not a conventional asymptotic ADM energy measured at spatial infinity. That distinction is essential in a closed universe.
\end{forthecurious}

\subsection{One Closure Principle, Two Applications}

We can now see the common structure without claiming that geometry and energy are literally the same equation.

For geometry, Premise~2 rejects an ungrounded boundary and an ungrounded global orientation distinction.

For global energy, Symmetric Selection rejects an ungrounded choice of sign and magnitude.

In both cases the same question is being asked: \textbf{What physical fact fixes this global choice?} If the answer is ``nothing,'' the choice is a Brute Fact and is excluded by the framework.

That is why the geometric closure and the zero-energy closure belong in the same philosophical story while remaining mathematically distinct.

\begin{plainwords}
The Whole is closed in shape and has a symmetry-selected global energy value. Local energies need not vanish. The zero-energy claim is a separate closure selection, not part of the $S^3$-bundle classification.
\end{plainwords}

\newpage
\section{Final Summary and the Branch Map}

We have now followed the argument from the three foundational premises to the geometric object that Document~II classifies and to the global energy selection. The purpose of this final section is to make the necessity structure visible in one place.

\subsection{The Necessity Ladder}

The argument is not a list of disconnected guesses. It is a sequence of forks. At each fork there are alternatives. The paper asks whether each alternative can preserve the same foundational requirements.

\begin{chainbox}
\begin{enumerate}
    \item \textbf{Premise~1 (Identity):} A determinate existent is itself, $A=A$.
    \item \textbf{Premise~2 (No Brute Facts):} No fundamental physical feature may simply be ``just because.''
    \item \textbf{Premise~3 (Physical Actualization):} A physical fact must be physically instantiated.
    \item \textbf{Existence:} Absolute Nothingness cannot be the actual physical state, so some determinate physical actuality exists.
    \item \textbf{Maximal Whole:} Compositional Realism treats the totality of physical existents as one real Whole $W$.
    \item \textbf{No Outside:} No physical thing can exist outside the totality $W$.
    \item \textbf{Connected Whole / No True Separation:} The physical domain is one connected totality.
    \item \textbf{Global Closure:} The Read-Head postulate requires the Whole to realize its own physical closure.
    \item \textbf{Finite Completion:} Finite Actualization excludes actual physical infinity as a fundamental physical content.
    \item \textbf{Temporal Minimality:} One temporal dimension with finite, connected, boundaryless closure gives $\Sone$.
    \item \textbf{Dimensional Minimality:} The framework requires three spatial dimensions and no additional ungrounded dimension.
    \item \textbf{Spatial Topology:} Compactness, boundarylessness, and simple connectivity leave $\Sthree$.
    \item \textbf{Bundle Classification:} Document~II shows that the stated $S^3$-bundle over $S^1$ has exactly two smooth bundle classes.
    \item \textbf{Handedness Elimination:} No Ungrounded Global Handedness excludes the orientation-preserving branch unless independently grounded.
    \item \textbf{Klein Block:} The orientation-reversing class is therefore the unique surviving bundle class, $\Klein$.
    \item \textbf{Zero-Energy Selection:} Symmetric Selection fixes the global energy value at the unique fixed point of $E \mapsto -E$, namely $E_{\mathrm{global}}=0$.
\end{enumerate}
\end{chainbox}

The decisive phrase in this ladder is \textbf{surviving branch}. The paper does not need to show that every alternative is unimaginable. It needs to show that every alternative either violates a foundational requirement or adds a new physical fact that the framework has no ground to accept.

\newpage
\subsection{The Five Closure-Admissible Postulates}

\cref{tab:postulates} collects the five statements this paper hands to Document~II~\cite{CanonWhat}. They are the bridge between the philosophy of Document~I and the mathematics of Document~II.

\begin{table}[h]
\centering
\small
\begin{tabularx}{\textwidth}{@{}>{\raggedright\arraybackslash}p{4.8cm}X p{1.7cm}@{}}
\toprule
\textbf{Postulate (Document~II)} & \textbf{Established Here} & \textbf{Section} \\
\midrule
II-2.1---Global Temporal Fibration & Global closure is represented by a smooth fibration over a connected, compact, one-dimensional temporal base with connected spacelike three-dimensional fibers. & 4.4, 8.1 \\
II-2.2---Spatial Compactness \& Boundarylessness & Spatial extension is finite and completed without a physical wall or edge. & 5, 7.3 \\
II-2.3---Spatial Simple Connectivity & Nontrivial spatial holes have no independent ground under the Topological Razor. & 7.3 \\
II-2.4---Time-Orientability & The ordered temporal role requires a consistent local future direction. & 6.4 \\
II-2.5---Nontrivial Vertical Orientation Monodromy & The orientation-preserving branch carries an ungrounded global fiber-orientation distinction and is excluded unless independently grounded. & 8.3--8.4 \\
\bottomrule
\end{tabularx}
\caption{The five closure-admissible geometric conditions passed from Document~I to Document~II.}
\label{tab:postulates}
\end{table}

These five statements are not imported as convenient assumptions after the fact. The purpose of the preceding sections is to show why the framework adopts each one as a required bridge condition if it is to preserve its foundational rules. Document~II then removes the remaining mathematical ambiguity.

\subsubsection*{Postulate Independence Audit}
For the Upstream Pressure Map (Appendix) to function correctly, the bridge conditions must be audited for logical independence. 
\begin{itemize}
    \item \textbf{Postulates 2.1, 2.2, 2.3, and 2.5 are mathematically independent.} One can construct smooth fibrations over $S^1$ that violate any one of these conditions while satisfying the other three (e.g., non-compact fibers violate 2.2; toroidal fibers violate 2.3; the trivial bundle violates 2.5).
    \item \textbf{Postulate 2.4 (Time-Orientability) is mathematically redundant.} As proven in Document~II (Remark 2.2), once a Lorentzian metric with spacelike fibers over the orientable base $S^1$ is established (Postulate 2.1), time-orientability follows automatically. 
\end{itemize}
\textbf{Consequence for the Pressure Map:} Postulate 2.4 is retained in the ledger because it represents a distinct \emph{physical architectural requirement} (the existence of a local causal arrow), but it does not act as an independent topological selector. Any physical failure of time-orientability routes upstream not to a distinct topological postulate, but to the existence of the Lorentzian fibration itself (Postulate 2.1) or the dynamical realization of the metric.

\subsection{The Branch Map}

This paper does not ask the reader to accept the conclusion merely because it has been stated. It issues a precise challenge:

\begin{keyresult}
\textbf{No-Escape Test:} To defeat the conclusion within the same ontology, exhibit a physical model that preserves all three foundational premises, satisfies the same closure requirements, and nevertheless produces a different topology or global energy value. If the proposed alternative instead rejects a premise or adds an ungrounded physical feature, it is a different ontology rather than a surviving branch of the present argument.
\end{keyresult}

The following Branch Map records the cost of the principal escape routes.

\begin{table}[h]
\centering
\small
\begin{tabularx}{\textwidth}{@{}>{\raggedright\arraybackslash}p{4.2cm} >{\raggedright\arraybackslash}p{5.2cm} X@{}}
\toprule
\textbf{If you reject\ldots} & \textbf{Then you permit\ldots} & \textbf{Cost to the framework} \\
\midrule
Premise~1 & No determinate identity & Existence becomes self-contradictory. The rejection is self-refuting. \\
Premise~2 & Fundamental brute facts & Explanation may terminate in ``just because.'' \\
Premise~3 & Uninstantiated ``physical'' truths & Physical actuality loses its truthmaker requirement. \\
Compositional Realism & No real Whole & Global closure loses its physical bearer. \\
Connected Whole / No True Separation & Disconnected physical islands & A single connected physical domain is abandoned. \\
Read-Head / Global Closure & No internal physical closure & The Whole remains abstract or externally completed. \\
Finite Actualization & Actual physical infinity & Finite determinate completion is abandoned. \\
Temporal Minimality & Extra temporal dimensions & Unsupported temporal structure enters. \\
Dimensional Minimality & Extra spatial dimensions & Unsupported spatial structure enters. \\
Spatial Simple Connectivity & Fundamental holes & Extra topology is accepted without ground. \\
No Ungrounded Global Handedness & Orientation-preserving bundle & Permanent global handedness is accepted without ground. \\
Symmetric Selection & Nonzero global energy & Global sign and magnitude become brute. \\
\bottomrule
\end{tabularx}
\caption{The Branch Map: the principal alternatives and the exact commitment each one must surrender or supplement. The table is not intended to declare every rejected worldview irrational. Its function is narrower and more exact: it makes the \emph{price of escape} explicit.}
\label{tab:branchmap}
\end{table}

\subsection{What Would Count Against It}

A genuine counterargument must attack the logic at a load-bearing point. The relevant possibilities are clear.

The framework would be pressured by a robust model requiring a true first moment, an absolute temporal endpoint, an actual physical infinity, an external physical creator, or a fundamentally open linear time. Local curvature measurements alone are not direct falsifiers of topology unless tied to a global model incompatible with finite closed simply connected spatial topology.

\textbf{Conceptual defeat:} show that one of the three foundational premises is incoherent, mutually inconsistent, or incapable of doing the explanatory work claimed for it.

\textbf{Structural defeat:} preserve all three premises but demonstrate that a different closure structure satisfies them without introducing an ungrounded feature. In particular, a successful rival would need to defeat one of the bridge steps to compact temporal closure, three-dimensional spatial extension, simple connectivity, or the handedness elimination.

\textbf{Mathematical defeat:} show that the hypotheses handed to Document~II do not imply the stated bundle classification, or produce a counterexample within the exact mathematical category used there.

\textbf{Physical defeat:} provide a well-supported physical realization that requires a structure excluded by the framework---for example, a physically necessary temporal endpoint, an actually infinite completed physical extent, a physical external context, or a different globally grounded bundle class.

A merely imaginable alternative is therefore not enough. A counterexample must survive the same rules.

\newpage
\subsection{Not Claimed List}

To make the boundary of the argument equally clear, the framework explicitly does not claim the following:

\begin{itemize}
    \item No Brute Facts is not claimed to be the only possible metaphysics. It claims that the Klein Block follows when that premise is retained.
    \item Hilbert's Hotel does not prove mathematical infinity inconsistent. It is used only to illustrate why mathematical cardinality is not the same as finite physical amount.
    \item Read-Head is not derived from Identity alone. Global physical self-closure is an explicit postulate of the framework.
    \item Three dimensions are not forced by pure logic without the extension requirement.
    \item Topology alone does not determine metric or full dynamics.
    \item Closed time does not automatically solve causal paradoxes.
    \item A compact universe is not claimed to have ordinary asymptotic ADM energy. The canonical Hamiltonian statement is a separate conditional result.
    \item The sign-reversal symmetry is not claimed to be a theorem of all existing physics. It is the symmetry used by the framework's global-energy selection principle.
    \item Metric scale, constants, matter content, and complete dynamical law are not derived.
    \item Born rule and quantum measurement are not completed.
    \item Weak parity violation and baryogenesis are not derived from topology alone.
    \item It does not claim the Klein Block topology alone generates a nontrivial pure-gravitational anomaly. Document~III proves the pure Pin$^+$ bordism class is zero; any nontrivial anomaly requires additional gauge or fermionic data.
    \item Local energies are not claimed to vanish.
    \item Time is not claimed to be a machine, repeated timeline, or sequence observed from meta-time.
\end{itemize}

\begin{plainwords}
The result is a conditional necessity argument. Reality cannot be Absolute Nothingness within the foundational ontology. Once the three premises and the closure rules are retained, the mathematically possible alternatives are progressively eliminated as admissible physical realizations. Document~II then closes the remaining topological fork and identifies the unique surviving smooth bundle class as the Klein Block. Zero energy is selected by the separate symmetry argument.
\end{plainwords}

\newpage
\section*{Appendix: Rival-Preservation Table}
\phantomsection
\addcontentsline{toc}{section}{Appendix: Rival-Preservation Table}

The framework does not claim to refute rival ontologies from the inside. Instead, it executes the No-Escape Principle by showing what each rival must accept once it rejects a core commitment. This table measures the cost of escape \emph{as evaluated from within the framework's Principle of Sufficient Reason}. 

A holder of a rival view (e.g., a Humean) may defend their base-level stopping point as a virtue of parsimony rather than a cost. The table does not attempt to steelman or refute those internal defenses here; it merely makes the structural trade-offs explicit. A full comparative evaluation, engaging each rival's best self-description and demonstrating the specific explanatory work they cannot perform, is deferred to the proposed companion study: \textit{Why No Brute Facts?}

\begin{center}
\small
\begin{tabularx}{\textwidth}{@{}>{\raggedright\arraybackslash}p{3.2cm} >{\raggedright\arraybackslash}p{3.8cm} X@{}}
\toprule
\textbf{Commitment rejected} & \textbf{Strongest known rival} & \textbf{Explanatory burden the rival must carry} \\
\midrule
No Brute Facts & Humean regularity / brute-contingency metaphysics & Fundamental facts terminate in ``just because.'' Explanation is abandoned at the base level. \\
Physical Actualization & Mathematical possibilism / uninstantiated structures & Physical reality becomes disconnected from physical instantiation; truthmakers are no longer required. \\
Compositional Realism & Mereological nihilism / parts-only ontology & No single physical bearer exists for global closure or global energy. \\
Read-Head / Global Closure & Open-universe or externally completed cosmology & Closure is either permanently incomplete or borrowed from an external context (violating No Outside). \\
Finite Actualization & Actual-infinite cosmology & The physical Whole lacks a determinate, completed finite total content. \\
Temporal Minimality & Two-time or multi-time frameworks & Extra temporal structure enters without being required by the sequencing function. \\
Dimensional Minimality & Higher-dimensional theories (e.g., string theory) & Extra dimensions require independent grounding, compactification mechanisms, and an explanation for the observed 3D macroscopic arena. \\
Spatial Simple Connectivity & Torus / lens-space cosmologies & Extra holonomy and winding structure (nontrivial $\pi_1$) are accepted as fundamental without independent ground. \\
No Ungrounded Global Handedness & Orientation-preserving bundle & A permanent global fiber-orientation structure is admitted without an independent physical ground. \\
Symmetric Selection & Nonzero global energy cosmologies & A global sign and magnitude must be selected or grounded by an additional, unexplained principle. \\
\bottomrule
\end{tabularx}
\end{center}

The framework does not attempt here to prove that No Brute Facts is the only rational metaphysics. It declares No Brute Facts as a foundational commitment and exposes the structural cost of rejecting it. Having mapped the price of escape, the ontological engine is complete; the mathematical uniqueness proof now begins in Document~II.

\newpage
\section*{Appendix: Upstream Pressure Map}
\phantomsection
\addcontentsline{toc}{section}{Appendix: Upstream Pressure Map}

Document~I cannot be falsified by a single physical experiment, because it is a framework constitution, not a physical model. However, the physical gates formulated in Document~III exert \emph{upstream pressure} on these philosophical commitments. A persistent failure across all admissible physical repairs pressures the upstream postulates that mandated the topology.

\begin{center}
\small
\begin{tabularx}{\textwidth}{@{}>{\raggedright\arraybackslash}p{4.2cm} >{\raggedright\arraybackslash}p{4.2cm} X@{}}
\toprule
\textbf{Document~III Gate Failure} & \textbf{Immediate Consequence} & \textbf{Upstream Pressure on Document~I} \\
\midrule
No twisted Einstein--matter fixed point & No nonsingular cyclic solution in the specified sector & Pressures the \textbf{Read-Head} postulate as a physically realizable closure relation. \\
No quotient-compatible QFT state & Quantum completion fails on the chronology-violating quotient & Pressures \textbf{Physical Actualization} of quantum fields on the selected topology. \\
Entropy gates fail collectively & No non-circular thermodynamic completion & Pressures \textbf{Temporal Minimality} and \textbf{Global Closure} if cyclic time is shown to be physically unviable for thermodynamics. \\
Full SM Lagrangian descent fails & The topology cannot host the intended Standard Model & Pressures \textbf{Dimensional Minimality}, \textbf{Spatial Simple Connectivity}, or the \textbf{Klein Block} selection if the required gauge/fermion structures demand a different topology or bundle class. \\
Mixed bordism invariant is trivial & The proposed anomaly route is trivial for this topology & Pressures the specific \textbf{Klein Block} selection if the required physical asymmetries cannot be supported by the orientation-reversing monodromy. \\
\bottomrule
\end{tabularx}
\end{center}

\begin{plainwords}
A gate failure first constrains the dynamical layer (Document~III). It becomes pressure on a Document~I postulate only if every admissible repair inside the same postulate class fails. The philosophical premises are not falsified directly; they are pressured by the persistent inability of the physical universe to inhabit the geometry they demand.
\end{plainwords}

\newpage
\section*{Appendix: The Complete Deductive Chain}
\phantomsection
\addcontentsline{toc}{section}{Appendix: The Complete Deductive Chain}

For the rigorous reader, the entire architecture is presented here in compressed form. The labels matter. A \emph{Premise} is a foundational rule. A \emph{Commitment} states the ontology being used. A \emph{Selection} applies the Topological Razor to a fork. A \emph{Mathematical Fact} is supplied by the classification cited from Document~II or standard mathematics. An \emph{Interpretation} is not load-bearing.

\textbf{PHASE I: LOGICAL FLOOR AND FOUNDATIONAL PREMISES}
\begin{enumerate}
    \item \textbf{[Premise]} Identity: every determinate existent is self-identical, $A=A$.
    \item \textbf{[Logical background]} Logical non-contradiction: $A$ and $\neg A$ cannot both hold in the same respect.
    \item \textbf{[Premise]} No Brute Facts: every fundamental feature has a sufficient reason.
    \item \textbf{[Premise]} Physical Actualization: every physical fact is physically instantiated.
\end{enumerate}

\textbf{PHASE II: EXISTENCE AND THE WHOLE}
\begin{enumerate}[resume]
    \item \textbf{[Consequence]} Absolute Nothingness is not an admissible actual physical state.
    \item \textbf{[Necessity]} Some determinate physical actuality therefore exists.
    \item \textbf{[Commitment]} Compositional Realism: the totality of physical existents is one real Whole $W$.
    \item \textbf{[Consequence]} No physical thing exists outside $W$.
    \item \textbf{[Commitment]} Connected Whole / No True Separation: the physical domain is one connected totality.
    \item \textbf{[Consequence]} The physical Whole is uncreated by any external physical creator.
\end{enumerate}

\textbf{PHASE III: THE RAZOR AND GLOBAL CLOSURE}
\begin{enumerate}[resume]
    \item \textbf{[Method]} Topological Razor: fundamental physical structure requires a ground.
    \item \textbf{[Principle]} No-Escape Principle: a genuine rival must preserve the same foundational requirements and survive every later cut.
    \item \textbf{[Postulate]} Read-Head / Global Self-Closure.
    \item \textbf{[Requirement]} Read-Head requires sequence, extension, and completion.
\end{enumerate}

\textbf{PHASE IV: FINITENESS AND TIME}
\begin{enumerate}[resume]
    \item \textbf{[Selection]} Finite Actualization.
    \item \textbf{[Consequence]} Actual physical infinity is excluded within the framework.
    \item \textbf{[Selection]} Boundarylessness: a special physical wall is not admitted without independent ground.
    \item \textbf{[Selection]} Temporal Minimality: one temporal dimension is retained.
    \item \textbf{[Consequence]} Connected compact one-dimensional boundaryless temporal closure gives $\Sone$.
    \item \textbf{[Selection]} Time-Orientability: a consistent local future direction is required.
    \item \textbf{[Dynamical Obligation]} Closed temporal topology must reconcile causal and thermodynamic behavior with global closure.
\end{enumerate}

\textbf{PHASE V: SPACE AND TOPOLOGY}
\begin{enumerate}[resume]
    \item \textbf{[Selection]} Dimensional Minimality: three spatial dimensions are the minimum sufficient choice under the stated extension requirement.
    \item \textbf{[Postulate]} Spatial Compactness and Boundarylessness.
    \item \textbf{[Postulate]} Spatial Simple Connectivity.
    \item \textbf{[Mathematical Fact]} Closed simply connected 3-manifolds are $\Sthree$; the smooth dimension-three classification is supplied by Moise~\cite{Perelman,Moise}.
\end{enumerate}

\textbf{PHASE VI: THE TWIST}
\begin{enumerate}[resume]
    \item \textbf{[Postulate]} Global Temporal Fibration.
    \item \textbf{[Mathematical Fact]} Document~II gives exactly two smooth $S^3$-bundle classes over $S^1$ in the stated category~\cite{CanonWhat}.
    \item \textbf{[Method]} Structural Burden Rule: positive global structure (global orientability) requires an independent ground.
    \item \textbf{[Selection]} No Ungrounded Global Handedness rejects the orientation-preserving (globally orientable) branch unless independently grounded.
    \item \textbf{[Mathematical Fact]} The orientation-reversing class is the unique surviving class, the Klein Block $\Klein$~\cite{CanonWhat}.
    \item \textbf{[Dynamical Obligation]} The topology must support the required field, causal, and anomaly structures in Document~III.
\end{enumerate}

\textbf{PHASE VII: GLOBAL ENERGY}
\begin{enumerate}[resume]
    \item \textbf{[Selection]} Symmetric Selection adopts the sign-reversal involution $E \mapsto -E$ for the global energy datum.
    \item \textbf{[Logical Consequence]} The unique fixed point is $E_{\mathrm{global}} = 0$.
    \item \textbf{[Conditional Physical Support]} On the lifted compact boundaryless $\Sthree$ slices, the canonical Hamiltonian satisfies $H_{\mathrm{canonical}} \approx 0$ on the constraint surface~\cite{Wald,ADM,CanonWhat}.
    \item \textbf{[Interpretation]} The canonical constraint and the framework's global zero-energy selection are distinct but compatible statements.
\end{enumerate}

\section*{Appendix: Adversarial Tests of the Necessity Claim}
\phantomsection
\addcontentsline{toc}{section}{Appendix: Adversarial Tests of the Necessity Claim}

The following objections are the strongest forms a hostile reader is likely to raise. Each reply identifies exactly what would have to be shown for the objection to defeat the argument.

\begin{enumerate}
    \item \textbf{``The Principle of Sufficient Reason is optional.''} \\
    \textbf{Reply:} Yes. The paper does not claim otherwise. But the central result is a necessity result \emph{under} No Brute Facts. Rejecting it is a direct rejection of Premise~2, not a counterexample that preserves the framework. The philosophical cost is explicit: a fundamental fact may terminate explanation.

    \item \textbf{``Nothingness is possible.''} \\
    \textbf{Reply:} Then defend Absolute Nothingness as an actual physical state while preserving Physical Actualization. The framework's objection is that Absolute Nothingness contains no physical actuality and therefore cannot be the physically actual state described by the term ``physical reality.'' The issue is not whether one can utter the word ``nothing''; it is whether an absence of all physical actuality qualifies as an actual physical state.

    \item \textbf{``Reality could be created by something outside the universe.''} \\
    \textbf{Reply:} If the creator is physical, it belongs to $W$ by definition and is not outside it. If the creator is nonphysical, the proposal changes the ontology by adding a nonphysical explanatory domain. That may be a different metaphysics, but it is not an external physical component of the Whole. The present paper therefore establishes uncreatedness of the \emph{physical} Whole.

    \item \textbf{``There need not be a physical Whole.''} \\
    \textbf{Reply:} That is a rejection of Compositional Realism. The paper states that commitment openly because the closure argument is about the physical totality itself. A mere collection of parts can be defended, but then the single physical bearer of global closure is removed.

    \item \textbf{``The Read-Head is simply assumed.''} \\
    \textbf{Reply:} Correct: global self-closure is an explicit postulate. The argument does not conceal this. The philosophical task is to show why a no-external-context physical Whole requires an internally realized closure relation. The mathematical task begins only after that postulate is stated.

\newpage

    \item \textbf{``Actual physical infinity is possible.''} \\
    \textbf{Reply:} Mathematical infinity is not disputed. A smooth manifold may also contain infinitely many coordinate points while representing a finite physical volume; Finite Actualization concerns physical content, not mathematical coordinates. Actual physical infinity is rejected by Finite Actualization, an explicit selection principle. Hilbert's Hotel illustrates the difference between infinite cardinality and ordinary finite physical amount; it is not itself the proof of physical finiteness.

    \item \textbf{``A finite universe can have a boundary.''} \\
    \textbf{Reply:} Of course. The framework excludes the boundary because a physically distinguished stopping surface is an additional structural fact. Under No Brute Facts, that fact requires a ground. No such ground is supplied by the framework.

    \item \textbf{``Why one time dimension?''} \\
    \textbf{Reply:} Because the Read-Head's sequencing function requires one ordered temporal dimension. A second is not needed for that function. To defeat Temporal Minimality, show that the physical closure requires a second temporal dimension or supply an independent physical ground for it.

    \item \textbf{``Why three space dimensions?''} \\
    \textbf{Reply:} Under the stated extension requirement, three dimensions are the minimum sufficient setting for general finite networks modeled as non-intersecting channels. To defeat the step, reject that requirement, realize the Read-Head's function in fewer dimensions, or accept extra dimensions while supplying: a physical ground, an exact-number selector, and an account of the observed three-dimensional arena. Without these, extra dimensions are ungrounded structure.

    \item \textbf{``The torus is just as legitimate as $S^3$.''} \\
    \textbf{Reply:} Mathematically, many spaces are legitimate. Document~II records that relaxing simple connectivity admits fibers with nontrivial $\pi_1$, such as tori and lens spaces. The question is which topology is admissible under the declared structural requirements. The torus contains nontrivial loop structure; Spatial Simple Connectivity rejects that additional topology unless an independent physical ground is supplied.

    \item \textbf{``The straight bundle is mathematically possible.''} \\
    \textbf{Reply:} Exactly. Document~II proves it. The philosophical question is not mathematical existence but admissibility. The orientation-preserving branch carries global fiber-orientation structure around the temporal loop. If that structure is fundamental, it requires an independent ground. Otherwise No Ungrounded Global Handedness rejects the branch.

    \item \textbf{``Orientation-preserving does not mean someone chose left rather than right.''} \\
    \textbf{Reply:} Agreed. The relevant claim is not that an agent chose a label. It is that the bundle permits a globally preserved fiber orientation as a structural fact. The framework treats that global distinction itself as something that requires a ground.

    \item \textbf{``Closed time produces paradoxes.''} \\
    \textbf{Reply:} The framework does not dissolve the causal issue by changing its name. Document~II proves that, in a Lorentzian realization with spacelike $S^3$ fibers, every smooth future-directed timelike vector field transverse to the fibers has a closed orbit, and no global real-valued time function exists~\cite{CanonWhat}. This is a genuine causal cost of the topology~\cite{HawkingChronology}. Document~III treats the quotient-compatible state and physical viability of such chronology-violating structure as explicit open tests~\cite{CanonHow}. The framework therefore accepts this as a dynamical obligation.

    \item \textbf{``Entropy cannot increase around a circle.''} \\
    \textbf{Reply:} Correct for a single-valued entropy function that is strictly monotone around the entire circle. Document~II proves the corresponding no-go result for a nonconstant scalar entropy potential factoring through the temporal base~\cite{CanonWhat}. Document~III further shows that a globally defined, quotient-compatible entropy current with pointwise nonnegative divergence would have vanishing integrated entropy production, and it formulates the Tolman entropy problem as an explicit open test~\cite{CanonHow}. The framework does not claim that the thermodynamic mechanism has been derived; any proposed reset or reversal must pass non-circularity and hidden-entropy tests.
    
    \item \textbf{``Zero energy is arbitrary.''} \\
    \textbf{Reply:} Not once Symmetric Selection is accepted. The argument is conditional on the global sign-reversal involution $E \mapsto -E$. Zero is its unique fixed point. Zero energy is also a separate closure selection, not one of the five postulates handed to Document~II; the bundle classification can be completed without it. To keep a nonzero value within the framework, provide an independent ground for the preferred sign and magnitude or reject the symmetry principle.

    \item \textbf{``The canonical Hamiltonian is not the same as total energy.''} \\
    \textbf{Reply:} Correct. The paper states the distinction explicitly. Document~II formulates the constrained Hamiltonian result on the infinite cyclic cover $S^3 \times \mathbb{R}$, where the lifted $S^3$ leaves are compact and boundaryless~\cite{CanonWhat}. The statement $H_{\mathrm{canonical}} \approx 0$ is not an ADM energy at spatial infinity and is not by itself the definition of an observable total energy for the compact quotient.

    \item \textbf{``Consciousness is not actually explained.''} \\
    \textbf{Reply:} Correct as a criticism of a completed consciousness theory. Consciousness and quantum measurement are interpretive extensions of the closure ontology, not load-bearing premises for the topology. Document~III explicitly does not complete the Born rule or quantum measurement~\cite{CanonHow}. The topological derivation survives if the interpretation of consciousness as physical self-reference is rejected.
\end{enumerate}

\section*{Appendix: Glossary}
\phantomsection
\addcontentsline{toc}{section}{Appendix: Glossary}

\begin{itemize}
    \item \textbf{$A=A$:} Identity. A determinate thing is itself.
    \item \textbf{Brute Fact:} A fundamental feature with no sufficient reason.
    \item \textbf{Sufficient Reason:} A ground or explanation that makes a fundamental fact intelligible within the chosen ontology.
    \item \textbf{Physical Actualization:} A physical fact must be physically instantiated.
    \item \textbf{Maximal Whole ($W$):} The totality of physical existents treated as one real physical Whole.
    \item \textbf{Compositional Realism:} The totality of physical existents is treated as one real, physical Whole.
    \item \textbf{Connected Whole:} The physical domain is one connected totality.
    \item \textbf{No True Separation:} The framework rejects gaps understood as regions of absolute non-being.
    \item \textbf{Topological Razor:} No additional fundamental physical structure without an independent ground.
    \item \textbf{Structural Burden Rule:} Positive global structure (such as global orientability) requires an independent ground; the absence of such structure is the default.
    \item \textbf{No-Escape Principle:} A genuine rival must preserve the same foundational requirements and survive every later elimination step.
    \item \textbf{Read-Head:} The internal physical closure relation of the Whole.
    \item \textbf{Finite Actualization:} The Whole is required to have finite, determinate physical content.
    \item \textbf{Temporal Minimality:} One temporal dimension is admitted because it supplies the required ordering function.
    \item \textbf{Dimensional Minimality:} Three spatial dimensions are the minimum sufficient choice under the stated extension requirement.
    \item \textbf{No Ungrounded Global Handedness:} A globally preserved fiber orientation is not admitted as fundamental without a ground.
    \item \textbf{Symmetric Selection:} A sign-neutral global quantity is fixed at the unique point invariant under the selected sign-reversal symmetry.
    \item \textbf{Klein Block ($\Klein$):} The orientation-reversing smooth $S^3$-bundle over $S^1$ selected by the framework.
    \item \textbf{Seam:} The endpoint identification by which the temporal fundamental domain closes.
    \item \textbf{Fibration:} A smooth organization of spacelike fibers over the temporal base.
    \item \textbf{Monodromy:} The fiber map induced after one circuit around the temporal base.
    \item \textbf{Time-Orientability:} A continuous local choice of which timelike direction is future.
    \item \textbf{Simply Connected:} Every closed loop can be continuously shrunk to a point.
    \item \textbf{Boundaryless:} The manifold has no manifold boundary.
    \item \textbf{Compact:} A topological property expressed by finite-subcover compactness; in this framework it represents completed finite extent, not finite cardinality.
    \item \textbf{Actual Physical Infinity:} An actually infinite physical collection or extent treated as a completed physical fact.
\end{itemize}

\section*{Appendix: Related Work}
\phantomsection
\addcontentsline{toc}{section}{Appendix: Related Work}

These works are comparisons, not authorities. The present framework differs by combining closure, no brute facts, physical actualization, topological elimination, and the Klein Block.

\begin{itemize}
    \item \textbf{Leibniz and the Principle of Sufficient Reason:} Leibniz made sufficient reason and contradiction central principles of rationalist metaphysics. The present paper adopts a deliberately strong form of the sufficient-reason demand; contemporary discussion emphasizes that such a principle is powerful but controversial~\cite{SEPPSR}.
    \item \textbf{Metaphysical grounding:} Contemporary metaphysics distinguishes grounding from ordinary efficient causation and studies whether grounding is explanatory, constitutive, asymmetric, and well-founded~\cite{SEPGrounding}. Document~I uses ``sufficient reason'' in this broader grounding sense.
    \item \textbf{Cosmological arguments:} Classical and contemporary cosmological arguments often ask whether a contingent reality can terminate explanation in itself or in brute facts~\cite{CosmoArg}. The present paper's No-Escape Principle turns that general debate into a structural elimination procedure.
    \item \textbf{Einstein (1917):} Einstein's early relativistic cosmology explored a finite spatial universe. The present construction differs by making temporal closure and nontrivial bundle monodromy part of the ontology.
    \item \textbf{G\"odel (1949):} G\"odel showed that general relativity admits cosmological solutions with closed timelike curves. The present paper accepts the causal issue as a real test of its closed-time geometry rather than treating it as settled by philosophical closure.
    \item \textbf{Hawking's chronology-protection conjecture:} Hawking argued that quantum effects may prevent the formation of closed timelike curves~\cite{HawkingChronology}. The existence of such arguments makes the causal viability of a closed temporal base an important physical question for later work.
    \item \textbf{Hartle--Hawking (1983):} The no-boundary proposal seeks a boundary-free cosmological description within quantum cosmology. The present paper's no-boundary argument is instead derived from a closure ontology and a topological selection rule.
    \item \textbf{Structural realism:} Structural realist approaches emphasize relations and structure rather than ordinary object-centered pictures. The present paper similarly treats geometry and relations as fundamental, but adds a specific closure-based necessity argument.
\end{itemize}

\section*{Acknowledgments}

This research was conducted entirely independently, without institutional affiliation or external support. The Three-Paper Architecture separates the philosophical necessity argument, the differential-topological classification, and the physical dynamical tests.

The author used AI language models as computational co-auditors and structural stress-testers, including for symbolic calculation, structural review, and document preparation. The author retains full responsibility for the concepts, arguments, citations, and final claims.

\newpage
\phantomsection
\addcontentsline{toc}{section}{References}
\begin{thebibliography}{99}
\CanonBibliographySetup

\bibitem{SEPPSR}
K.~Easwaran, A.~H\'ajek, P.~Mancosu, and G.~Oppy, ``Principle of Sufficient Reason,'' \textit{The Stanford Encyclopedia of Philosophy}, substantive revision June 14, 2023, E.~N.~Zalta and U.~Nodelman (eds.), \url{https://plato.stanford.edu/entries/sufficient-reason/}.
\textit{Background on the Principle of Sufficient Reason, its formulations, and the distinction between causal and broader sufficient-reason accounts.}

\bibitem{SEPGrounding}
T.~A.~Raven, ``Metaphysical Grounding,'' \textit{The Stanford Encyclopedia of Philosophy}, substantive revision December 6, 2021, E.~N.~Zalta (ed.).
\textit{Background on grounding as a form of constitutive determination or explanation and on the relation between grounding and necessity.}

\bibitem{CosmoArg}
E.~N.~Zalta (ed.), ``Cosmological Argument,'' \textit{The Stanford Encyclopedia of Philosophy}.
\textit{Background on arguments from sufficient reason, contingency, and explanatory termination.}

\bibitem{Perelman}
G.~Perelman, ``The entropy formula for the Ricci flow and its geometric applications,'' \textit{arXiv:math/0211159} (2002); ``Ricci flow with surgery on three-manifolds,'' \textit{arXiv:math/0303109} (2003); ``Finite extinction time for the solutions to the Ricci flow on certain three-manifolds,'' \textit{arXiv:math/0307245} (2003).
\textit{Together these works contain the Ricci-flow results underlying the proof of the Poincar\'e conjecture used for the closed simply connected 3-manifold step.}

\bibitem{Moise}
E.~E.~Moise, \textit{Geometric Topology in Dimensions 2 and 3}, Graduate Texts in Mathematics, Vol.~47, Springer, 1977.
\textit{Provides the dimension-three smoothing and triangulation results used to pass from the topological $S^3$ classification to the smooth category.}

\bibitem{Klein}
F.~Klein, ``\"Uber Riemann's Theorie der algebraischen Functionen und ihrer Integrale,'' \textit{Mathematische Annalen} \textbf{21}, 141--162 (1882).
\textit{Historical source associated with Klein's work; the name ``Klein Block'' is used here by analogy with the classical Klein bottle.}

\bibitem{SEPInfinity}
K.~Easwaran, A.~H\'ajek, P.~Mancosu, and G.~Oppy, ``Infinity,'' \textit{The Stanford Encyclopedia of Philosophy} (Summer 2025 Edition), E.~N.~Zalta and U.~Nodelman (eds.), \url{https://plato.stanford.edu/archives/sum2025/entries/infinity/}.
\textit{Background for distinguishing mathematical infinity from claims about actual physical infinity, including the Hilbert Hotel discussion.}

\bibitem{Wald}
R.~M.~Wald, \textit{General Relativity}, University of Chicago Press, 1984.
\textit{Standard reference for the geometric and canonical formulations of General Relativity.}

\bibitem{ADM}
R.~Arnowitt, S.~Deser, and C.~W.~Misner, ``The Dynamics of General Relativity,'' in L.~Witten (ed.), \textit{Gravitation: An Introduction to Current Research}, Wiley, 1962, pp.~227--264; reprinted in \textit{General Relativity and Gravitation} \textbf{40}, 1997--2027 (2008).
\textit{Foundational reference for the ADM Hamiltonian formulation and its boundary terms.}

\bibitem{HawkingChronology}
S.~W.~Hawking, ``Chronology protection conjecture,'' \textit{Physical Review D} \textbf{46}, 603--611 (1992), \href{https://doi.org/10.1103/PhysRevD.46.603}{doi:10.1103/PhysRevD.46.603}.
\textit{Classic discussion of the physical difficulties associated with closed timelike curves and chronology protection.}

\bibitem{LeeYang}
T.~D.~Lee and C.~N.~Yang, ``Question of Parity Conservation in Weak Decays,'' \textit{Physical Review} \textbf{104}, 254--258 (1956).
\textit{Theoretical proposal questioning parity conservation in weak interactions.}

\bibitem{Wu}
C.~S.~Wu, E.~Ambler, R.~W.~Hayward, D.~D.~Hoppes, and R.~P.~Hudson, ``Experimental Test of Parity Conservation in Beta Decay,'' \textit{Physical Review} \textbf{105}, 1413--1415 (1957).
\textit{Experimental confirmation of parity violation in weak interactions.}

\bibitem{Hatcher}
A.~E.~Hatcher, ``A proof of the Smale Conjecture, $\Diff(S^3) \simeq O(4)$,'' \textit{Annals of Mathematics} \textbf{117}, 553--607 (1983).
\textit{Establishes that the diffeomorphism group of the 3-sphere deformation-retracts onto the orthogonal group, restricting the mapping class group to $\mathbb{Z}_2$.}

\bibitem{CanonWhat}
Canon, ``The Klein Block: Topological Uniqueness of the Non-Orientable $S^3$-Bundle over $S^1$,'' (Companion Mathematical Manuscript, Document~II of the Three-Paper Architecture), Zenodo, 2026, \href{https://doi.org/10.5281/zenodo.22766247}{doi:10.5281/zenodo.22766247}.

\bibitem{CanonHow}
Canon, ``Twisted Dynamics on the Klein Block: A Conditional Framework for Kinematic Descent, Nodal Defect Structure, and Anomaly Constraints,'' (Companion Physics Manuscript, Document~III of the Three-Paper Architecture), Zenodo, 2026, \href{https://doi.org/10.5281/zenodo.22766260}{doi:10.5281/zenodo.22766260}.

\end{thebibliography}

\end{document}